\documentclass[11pt,a4paper]{article}
\usepackage[margin=22mm]{geometry}
\usepackage{fontspec}
\setmainfont{Latin Modern Roman}
\setsansfont{Latin Modern Sans}
\setmonofont[Scale=0.85]{DejaVu Sans Mono}
\usepackage{babel}
\babelprovide[main,import]{czech}
\usepackage{amsmath,amssymb,amsthm,mathtools}
\renewcommand{\proofname}{Důkaz}
\usepackage{booktabs,longtable,array}
\usepackage{enumitem}
\usepackage{xcolor}
\usepackage{hyperref}
\usepackage{xurl}
\usepackage{mathrsfs}
\usepackage{fancyhdr}
\hypersetup{colorlinks=true,linkcolor=blue!45!black,urlcolor=blue!45!black,pdftitle={TWIST-J: číselná okolí skutečných rozptylových prahů},pdfauthor={Analytická pracovní poznámka}}
\setlength{\parindent}{0pt}
\setlength{\parskip}{5pt}
\setlength{\emergencystretch}{3em}
\setlist{nosep,leftmargin=1.6em}
\pagestyle{fancy}
\fancyhf{}
\fancyhead[L]{\small\sffamily TWIST-J / číselná okolí rozptylových prahů}
\fancyhead[R]{\small\sffamily NON-CANONICAL / L1}
\fancyfoot[C]{\small\thepage}
\setlength{\headheight}{15pt}
\newcommand{\norm}[1]{\left\lVert#1\right\rVert}
\newcommand{\abs}[1]{\left\lvert#1\right\rvert}
\newcommand{\ind}{\mathbf1}
\newcommand{\dd}{\,\mathrm d}
\newcommand{\Rea}{\operatorname{Re}}
\newcommand{\Ran}{\operatorname{Ran}}
\newcommand{\Ccal}{\mathcal C}
\newcommand{\Tcal}{\mathcal T}
\newcommand{\eps}{\varepsilon}
\newtheorem{theorem}{Tvrzení}[section]
\newtheorem{lemma}[theorem]{Pomocné tvrzení}
\theoremstyle{remark}\newtheorem{remark}[theorem]{Poznámka}
\numberwithin{equation}{section}
\providecommand{\tightlist}{\setlength{\itemsep}{0pt}\setlength{\parskip}{0pt}}

\begin{document}

\thispagestyle{empty}
{\Large\sffamily\bfseries Číselná okolí skutečných rozptylových prahů TWIST-J\par}
\vspace{3mm}
{\large Úplné místní periody a původní přenosový sloupec\par}
\vspace{4mm}
Samostatný analytický dodatek, 29. října 2026\par
\textbf{NON-CANONICAL / L1. Interně přezkoumaný návrh důkazu.}

\section*{Nový výsledek}

Pro všech osm setkání objektového pólu s kontinuální hranou,
oba dosažitelné kuželové konce a oba průsečíky výměny pulzních
větví s hranou jsou nyní doloženy společné výslovné poloměry
\[
\boxed{r=2^{-24},\qquad R_0=2^{-36},\qquad \rho=2^{-64}.}
\]
První číslo určuje pevný integrační obvod, druhé komplexní
analytické okolí jeho periody a třetí okolí přesné shody
s původním vnějším fyzikálním předpisem. Původní pólové,
kontinuální i pulzní příspěvky jsou zachovány. Všechna tato
místní pokračování mají výslovné meze libovolných derivací.

Pro oba skutečné stacionární vázané prahy je navíc odvozen
přesný kvantitativní normální tvar
\[
\boxed{F_{2,\mathrm{loc}}(K,E)
=\frac{A(K,Z)}{\sqrt Z}+B_{\mathrm{loc}}(K,Z),
\qquad Z=E-E_b(K),}
\]
\[
|K-K_0|<2^{-88},\quad |Z|<2^{-152},\qquad
2^{-8}<|A|<2^{12},\quad |B_{\mathrm{loc}}|<2^{80}.
\]
Pro reálné \(K\) a \(E\in I\), na spektrální straně \(Z>0\), z něj bez horní meze
komplementu plyne odhad \textbf{celého} původního přenosu
\[
\boxed{\left\|V(K,E)-\binom{-2/F_1(E)}0\right\|_2
\le 2^{20}\sqrt Z.}
\]
Celý první kontakt \(F_1\) již má rovněž vyčíslené komplexní
okolí s \(\Re F_1>1/16\). Pro oboustranný normální tvar
celého \(V\) zbývá omezit úplný komplement druhého kontaktu.
Tento rozdíl je v důkazu výslovný.

\vfill
\small
Dodatek neurčuje menší globální čas. Dřívější samostatný
časový návrh [S4] zůstává beze změny. Jde o nové místní
analytické výsledky pro původní normalizovaný vstup, pevný
počátek času a úplný vybraný jev. Všech 26 štítků a korelace
paměti s referencí se zachovávají. Práce je oddělena od
rezervace \#1816 a nepředstavuje uskutečněný experiment.
\newpage
\normalsize


\section{Původní úloha, normalizace a návaznost}

\subsection{Přesně zachovaný zdroj a jev}

Používáme původní unitární model, úplné kontaktní funkce
a rozptylové předpoklady z [S1--S3]. Neprovádí se nová
postselekce ani změna zdrojové normalizace. Položme
\[
h=10^{-3},\quad E_*=8\pi/5,\quad I=[E_*-h,E_*+h],
\quad \beta=4\pi/5,\quad e=E-\beta,\quad g=2\pi/5.
\]
Původní normalizovaný objektový vstup je
\[
\psi_{\rm obj}
=[2h\rho_{\rm in}(H_1)]^{-1/2}
\mathbf1_I(H_1)\mathbf1_{[-h,h]}(P)a_0,
\quad a_0=\delta_{R=0,r=0}\,iB,
\quad \|\psi_{\rm obj}\|=1.
\]
Token je přesně v \(y=0\). Celková hybnost \(K\) proto
zůstává rovnoměrná na celém kruhu; energetická míra vstupu
je přesně \(\mathbf1_I(E)\,\mathrm dE/(2h)\).
Výstupní projektor je stále
\[
\begin{aligned}
Q={}&\mathbf1_{I_K}(K)\mathbf1_{|P|\le1/100}
\mathbf1_I\bigl(\Omega_b(P)+\beta+\varepsilon(k)\bigr)\Pi_b(P)\\
&\otimes P_+(k)\mathbf1_{\varepsilon'(k)>0}
\otimes I_{\rm label}\otimes|d_1\rangle\langle d_1|
\otimes I_{\rm ref}.
\end{aligned}
\]
Operátor \(I_{\rm label}\) zachovává všech 26 štítků.
Původní paměť může být v libovolné korelaci s referencí.
Místní dekompozice integračního kruhu níže slouží důkazu
regularity; žádný člen zdroje ani jevu se tím nevynechává.

\subsection{Úplný kontakt a rozptylový operátor}

Původní pulzní fáze a váhy jsou
\[
\varepsilon(k)=\arccos\frac{1-4\cos k}{5},\qquad
h_\sigma(k)=\frac{1-\sigma\varepsilon'(k)}2.
\]
Se správným vnějším předpisem \(E\mapsto E+i0\) platí
\[
F_1(E)=\int_{-\pi}^{\pi}f_B(P,E-\beta)\frac{\mathrm dP}{2\pi},
\]
\[
F_2(K,E)=\sum_{\sigma=\pm1}\int_{-\pi}^{\pi}
h_\sigma(K-P)f_A\bigl(P,E-\beta-\sigma\varepsilon(K-P)\bigr)
\frac{\mathrm dP}{2\pi},
\]
\[
T(K,E)=\frac1{4+F_1F_2}
\begin{pmatrix}-2F_2&-4\\4&-2F_1\end{pmatrix},\qquad
V(K,E)=T(K,E)e_1=\frac{(-2F_2,4)^T}{4+F_1F_2}.
\]
V druhém kontaktu je ponecháno celé objektové kontinuum,
všechny kontaktně viditelné póly a obě pulzní větve. Symbol
\(V(k)\) v místní maticové identitě pro pulz níže označuje
původní dvourozměrnou pulzní matici. Je odlišný od sloupce
\(V(K,E)\); argumenty obou symbolů jsou vždy uvedeny.

\subsection{Přehled nových číselných mezí}

U tří regulárních typů místních příspěvků je společný
komplexní polydisk \(|K-K_0|,|e-e_0|<R_0=2^{-36}\).
Na reálném polydisku \(\rho=2^{-64}\) se prokazuje
shoda s původní vnější periodou.

\begin{center}
\begin{tabular}{p{65mm}cc}
\toprule
Místní příspěvek & Mez amplitudy & Značky důkazu\\
\midrule
Pól a kontinuální hrana společně & \(2^{-3}\) & QA1--QA23\\
Obě pulzní větve při výměně & \(2^{-17}\) & QB1--QB16\\
Obě hrany kuželového konce & \(2^{-18}\) & QC1--QC23\\
\bottomrule
\end{tabular}
\end{center}

Je-li příslušný exponent v tabulce \(c=3,17,18\), pak
ve středu platí
\[
|\partial_K^m\partial_E^n I|\le m!n!\,2^{36(m+n)-c}.
\]
Na fyzikálně identifikovaném okolí \(\rho\) se 36 nahradí 37.
Tato tabulka neomezuje celý \(F_2\); každá uvedená perioda
je přesně vymezenou částí jeho integrálu.

\subsection{Převzaté již vyčíslené středové údaje}

Úplný soupis fyzikálních pólů a stacionárních konfigurací
je v předchozím samostatném dodatku [S5]. Označení BS, CS,
PE odkazují na jeho rovnice. Označení PT odkazuje na
oddíl 14 předchozího časového návrhu [S4]. Zde se přebírá
zejména následující:
\begin{enumerate}
\item V původním okně jsou právě dvě stacionární vázané
konfigurace, obě na hlavní větvi posunuté o \(\pi\),
s pulzním znaménkem \(\sigma=-1\). Jejich společná
křivost je větší než \(19/64\), objektová veličina
splňuje \(|t_0|<2\), reziduum \(r_b>1/32\) a pulzní
váha \(h_->3/8\). Rovnice BS25--BS31 určují původní
vrstvu i znaménka.
\item Kruhová vzdálenost těchto vázaných křivek od obou
běžných kontinuálních stacionárních křivek je větší
než \(1/300\), podle CS12. Tento rozestup nenahrazuje
číselnou kontrolu zbytku celé periody.
\item Setkání objektového pólu a hrany nastává na osmi
křivkách PT25. Pro ně platí společné nenulové meze
uvedené přímo v QA6--QA7. Dosažitelný kužel a výměna
pulzních větví jsou určeny v PT28--PT31.
\end{enumerate}
Nové výsledky tohoto dodatku jsou číselná okolí a celé
místní analytické rozvoje. Nepřidávají nový fyzikální
předpoklad. Obecná existence vlnových operátorů a původní
model zůstávají převzatými vstupy [S1--S3].

\section{Pevný obvod pro všech osm setkání pólu s hranou}
\hypertarget{vuxfdsledek}{%
\subsection{Výsledek}\label{vuxfdsledek}}

Fixujme kterýkoli bod \((P_0,K_0,e_0)\) z osmi křivek PT25, kde
\(e_0\in I_e=[4\pi/5-1/1000,4\pi/5+1/1000]\). Zvolme

\[
\boxed{r=2^{-24},\qquad R_0=2^{-36},\qquad \rho=2^{-64}.}
\tag{QA1}
\]

Obvod \(\Gamma\) je půlkruh \(|P-P_0|=r\), vedený od \(P_0-r\) do
\(P_0+r\) v polorovině opačné ke znaménku
\(\Delta v=\omega_{\rm edge}'(P_0)-\sigma\varepsilon'(K_0-P_0)\).
Pokračujeme původním vnějším objektovým listem.

Na komplexním polydisku \(|K-K_0|,|e-e_0|<R_0\) je místní perioda po
tomto pokračování holomorfní a

\[
\boxed{
I_\Gamma(K,e)=\int_\Gamma h_\sigma(K-P)
f_A(P,e-\sigma\varepsilon(K-P))\frac{\mathrm dP}{2\pi},
\qquad |I_\Gamma|<257\sqrt r<\frac18.
}
\tag{QA2}
\]

Pro reálné \(|K-K_0|,|e-e_0|<\rho\) je to přesně původní vnější
příspěvek intervalu \([P_0-r,P_0+r]\), včetně fyzikálního pólového
členu, pokud je přítomen. Reziduum se neopravuje samostatně.

Ve středu tedy pro \(m,n\ge0\) platí

\[
\boxed{
|\partial_K^m\partial_E^nI_\Gamma(K_0,e_0)|
\le m!n!\,2^{36(m+n)-3}.
}
\tag{QA3}
\]

Na celém fyzikálně identifikovaném okolí poloměru \(\rho\) platí

\[
\boxed{
|\partial_K^m\partial_E^nI_\Gamma(K,e)|
\le m!n!\,2^{37(m+n)-3}.
}
\tag{QA4}
\]

Zde \(\partial_E=\partial_e\). \(R_0\) je velikost analytického
pokračování na pevném obvodu; \(\rho\) je vyčíslené okolí shody s
původním fyzikálním předpisem.

\hypertarget{pux16fvodnuxed-funkce-a-stux159edovuxe9-meze}{%
\subsection{Původní funkce a středové
meze}\label{pux16fvodnuxed-funkce-a-stux159edovuxe9-meze}}

Používáme přesně

\[
\omega=e-\sigma\varepsilon(k),\quad k=K-P,\quad v=\cos P,\quad
x=e^{-i(\omega-g)},\quad t=v-5\cos\omega,\quad g=2\pi/5,
\] \[
C=x^2-\frac25vx-\frac35,\quad
B=\frac35x^2+\frac25vx-1,\quad
q=\frac{t-\mathcal R}{4},\quad \mathcal R^2=t^2-16,
\] \[
D_o=xC+qB,\quad N=x^2+\frac35(v+q)x+vq,\quad
f_A=\frac{2xN}{D_o}-1,\quad h_\sigma=\frac{1-\sigma\varepsilon'(k)}2.
\tag{QA5}
\]

\(s=\pm1\) označuje objektovou hranu \(t=4s\). Objektová setkání jsou
\((s,v_0,\omega_0)=(-1,v_a,g/2)\) a \((1,-v_a,\pi+g/2)\), kde
\(v_a=5\cos(g/2)-4\). Předchozí přímá kontrola původních vzorců dává
jednotně

\[
|v_0|<1/8,\quad \frac9{16}<|t_{P,0}|<6,\quad
|B_0|=\frac45(\sqrt5-1)>\frac45,\quad D_{a,0}=0,
\tag{QA6}
\] \[
D_a=xC+\frac t4B,\qquad D_o=D_a-\frac B4\mathcal R.
\]

Dále \(|\Delta v|>9/40\), \(|\omega_{\rm edge}'|<3/8\) a
\(|\varepsilon'|<9/10\), takže

\[
\frac8{11}<\left|\frac1{\Delta v}\right|<\frac{40}{9}.
\tag{QA7}
\]

Meze platí na všech osmi křivkách a celém uzavřeném \(I_e\).

\hypertarget{vuxfdslovnuxe9-holomorfnuxed-pulznuxed-okoluxed}{%
\subsection{Výslovné holomorfní pulzní
okolí}\label{vuxfdslovnuxe9-holomorfnuxed-pulznuxed-okoluxed}}

V centrech obou rodin je \(\varepsilon_0=3\pi/5+\xi\) nebo
\(2\pi/5+\xi\), \(|\xi|\le1/1000\). Z \(\cos(2\pi/5)<31/100\) plyne
\(|\cos\varepsilon_0|<311/1000<1/3\).

Pišme \(a(k)=(1-4\cos k)/5=\cos\varepsilon(k)\). Pro \(|k-k_0|\le1/64\)
jsou \(|\sin k|,|\cos k|<2\), a tedy \[
|a-a_0|\le\frac85|k-k_0|\le\frac1{40},\qquad |a|<\frac12.
\] Funkce \(b(k)=\sqrt{1-a(k)^2}\) má jednoznačnou analytickou větev s
\(b(k_0)=\sin\varepsilon_0>0\), neboť \(|a^2|<1/4\). Její modul je větší
než \(4/5\). Původní pulzní fáze se holomorfně pokračuje a \[
\varepsilon'=-\frac{a'}b,\qquad
\varepsilon''=-\frac{a''}b-\frac{a(a')^2}{b^3}.
\] Z \(|a'|,|a''|<8/5\) dostáváme komplexní meze

\[
\boxed{|\varepsilon'|<2,\qquad |\varepsilon''|<5.}
\tag{QA8}
\]

Na pracovním boxu \(|P-P_0|,|K-K_0|,|e-e_0|\le2^{-10}\) je
\(|k-k_0|\le2^{-9}\), \(|\omega-\omega_0|\le5\cdot2^{-10}\), a proto
\(|x|<2\), \(|v|<1/4\), \(|\sin\omega|,|\cos\omega|<2\). Přímé derivace
dávají

\[
\boxed{|t_{PP}|<100,\qquad |t_K|<20,\qquad |t_e|<10.}
\tag{QA9}
\]

Například
\(t_{PP}=-\cos P+5\cos\omega\,\omega_P^2+ 5\sin\omega\,\omega_{PP}\) má
modul menší než \(2+40+50=92\).

\hypertarget{oddux11blenuxed-odmocninovuxfdch-nul}{%
\subsection{Oddělení odmocninových
nul}\label{oddux11blenuxed-odmocninovuxfdch-nul}}

Položme \(p=P-P_0\), \(u=t-4s\). Pro \(|p|=r\) a
\(|K-K_0|,|e-e_0|\le R_0\) dává (QA9)

\[
|u-t_{P,0}p|\le50r^2+30R_0<r/16.
\tag{QA10}
\]

Poslední porovnání je přesně \(50/2^{24}+30/2^{12}<1/16\). Z (QA6) proto
na celém kruhu

\[
\boxed{
r/2<|u|<7r,\quad 7<|t+4s|<9,\quad
\sqrt r<|\mathcal R|<8\sqrt r.
}
\tag{QA11}
\]

Poslední meze platí pro oba odmocninové listy. Rouchého věta pro
\(u=t_{P,0}p+\text{chyba}\) dává právě jeden jednoduchý kořen
\(P_e(K,e)\) uvnitř kruhu. Druhá hrana zde není.

\hypertarget{dolnuxed-mez-pux16fvodnuxedho-jmenovatele}{%
\subsection{Dolní mez původního
jmenovatele}\label{dolnuxed-mez-pux16fvodnuxedho-jmenovatele}}

Na témže obvodu

\[
|\omega-\omega_0|\le2r+3R_0<3r,\qquad
|x-x_0|<6r,\qquad |v-v_0|<2r.
\tag{QA12}
\]

Polynom \(B\) má na boxu \(|B_x|<5/2,\ |B_v|<4/5\), a tedy

\[
|B-B_0|<32r,\qquad \boxed{|B|>3/4.}
\tag{QA13}
\]

Pro nezávislé proměnné \(x,v,t\) s \(|x|<2,|v|<1/4,|t|<5\) má polynom
\(D_a\) \[
|(D_a)_x|<17,\qquad |(D_a)_v|<3,\qquad |(D_a)_t|<1.
\] Protože \(D_{a,0}=0\), platí

\[
\boxed{|D_a|<17(6r)+3(2r)+7r=115r<128r.}
\tag{QA14}
\]

Úplný původní jmenovatel na obou listech splňuje

\[
\boxed{
|D_o|\ge\frac{|B||\mathcal R|}{4}-|D_a|
>\frac3{16}\sqrt r-128r
=\frac5{32}\sqrt r>\frac18\sqrt r.
}
\tag{QA15}
\]

Rovnost používá přesnou volbu \(\sqrt r=2^{-12}\).

\hypertarget{uvnitux159-je-pruxe1vux11b-jeden-prux16fmux11bt-puxf3lu}{%
\subsection{Uvnitř je právě jeden průmět
pólu}\label{uvnitux159-je-pruxe1vux11b-jeden-prux16fmux11bt-puxf3lu}}

Eliminace odmocniny v \(D_o=0\) dává analytický polynom

\[
\mathcal L=D_a^2-\frac{B^2}{16}(t^2-16)
=(xC)^2+\frac t2(xC)B+B^2.
\tag{QA16}
\]

V celém disku je \(B\ne0\) díky (QA13). Na obvodu jsou \[
|D_a|^2<2^{14}r^2,\qquad
\left|\frac{B^2}{16}(t^2-16)\right|>\frac{63}{512}r.
\] První výraz je menší než druhý. Druhá Rouchého aplikace tedy dává
právě jeden jednoduchý kořen \(\mathcal L\), označený \(P_b(K,e)\). Je
to průmět pólu na některém z obou objektových listů; neprohlašuje se za
fyzikální atom na obou stranách setkání.

Jednoduchost a jednoznačnost dávají společnou holomorfnost obou kořenů
na celém parametrickém polydisku \(R_0\), s

\[
|P_e-P_0|<r,\qquad |P_b-P_0|<r.
\tag{QA17}
\]

Přímo lze použít konturový vzorec \[
P_b-P_0=\frac1{2\pi i}\oint_{|p|=r}
p\,\frac{\partial_p\mathcal L}{\mathcal L}\,\mathrm dp.
\]

\hypertarget{pux159uxedmuxe1-reuxe1lnost}{%
\subsubsection{Přímá reálnost}\label{pux159uxedmuxe1-reuxe1lnost}}

Pro \(\delta=\omega-g,\ c=\cos\delta,\ s_\delta=\sin\delta,\ A=c-v\)
platí \(C=(2/5)x(A-4is_\delta)\), \(B=-(2/5)x(A+4is_\delta)\). Dosazení
dává přesně

\[
\boxed{\mathcal L=\frac2{25}x^3\Phi(v,\omega),}
\tag{QA18}
\] \[
\Phi=(4c-v+5\cos\omega)(c-v)^2
-16s_\delta^2(6c-v-5\cos\omega).
\]

Pro reálné parametry je \(\Phi\) reálná analytická funkce. Jednoznačný
kořen v disku symetrickém vůči reálné ose je tedy reálný. Stejná
vlastnost \(P_e\) plyne z reálnosti \(t\). Není zde použit nevyčíslený
implicitní poloměr.

\hypertarget{vyux10duxedslenuxfd-spoleux10dnuxfd-vnux11bjux161uxed-smux11br}{%
\subsection{Vyčíslený společný vnější
směr}\label{vyux10duxedslenuxfd-spoleux10dnuxfd-vnux11bjux161uxed-smux11br}}

Ve středu je \(D_a=0,\ t=4s,\ P_e=P_b=P_0\). Derivace (QA16) proto
dávají

\[
\partial_eP_b=-\frac{\mathcal L_e}{\mathcal L_P}
=-\frac{t_e}{t_P}
=\partial_eP_e=\frac1{\Delta v}.
\tag{QA19}
\]

Z (QA17) a Cauchyových mezí na polovičním polydisku \(R_0/2\) plyne pro
obě \(P_j,\ j=e,b\),

\[
|\partial_KP_j|,\ |\partial_eP_j|\le2r/R_0=2^{13},
\] \[
|\partial_e\partial_KP_j|,\ |\partial_e^2P_j|
\le8r/R_0^2=2^{51}.
\tag{QA20}
\]

Na malém komplexním polydisku \(\rho=2^{-64}\) proto

\[
|P_j-P_0|\le2^{14}\rho=2^{-50}<r/4,\qquad
\left|\partial_eP_j-\frac1{\Delta v}\right|
\le2^{52}\rho=2^{-12}.
\tag{QA21}
\]

Podle (QA7) má reálná část obou energetických derivací stejné znaménko
jako \(\Delta v\) a modul větší než \(1/2\). Pro reálné parametry v
malém okolí a malé \(\epsilon>0\), které ponechá \((K,e+i\epsilon)\)
uvnitř téhož polydisku, integrace derivace po svislé úsečce dává

\[
\boxed{
\operatorname{sgn}\operatorname{Im}P_e(K,e+i\epsilon)
=\operatorname{sgn}\operatorname{Im}P_b(K,e+i\epsilon)
=\operatorname{sgn}\Delta v.
}
\tag{QA22}
\]

V disku nejsou jiné objektové body větvení, jiné průměty pólů ani pulzní
degenerace. Reálný interval lze tedy v opačné polorovině deformovat na
pevný půlkruh \(\Gamma\), bez překročení singularity. Po celou deformaci
se pokračuje původním vnějším listem. Limit \(\epsilon\downarrow0\)
zachovává přesně původní místní periodu, včetně fyzikálního rezidua tam,
kde pól existuje.

\hypertarget{uxfaplnuxe1-muxedstnuxed-perioda-a-derivace}{%
\subsection{Úplná místní perioda a
derivace}\label{uxfaplnuxe1-muxedstnuxed-perioda-a-derivace}}

Na obvodu z (QA11) plyne \(|q|<2\). Proto \[
|N|<4+\frac35(1/4+2)\,2+\frac14\,2<8,\qquad |h_\sigma|<2,
\] a podle (QA15) \[
|f_A|<256/\sqrt r+1.
\]

Na celém polydisku \(R_0\) jsou odmocninové nuly oddělené od obvodu a
jmenovatel nenuluje. List lze podél oblouku a v jednoduše souvislém
parametrickém polydisku pokračovat společně. Integrand a jeho perioda
jsou tedy holomorfní. Z délky oblouku \(\pi r\) plyne

\[
|I_\Gamma|
\le\frac{\pi r}{2\pi}\,2\left(\frac{256}{\sqrt r}+1\right)
=256\sqrt r+r<257\sqrt r<\frac18.
\tag{QA23}
\]

Cauchyův vzorec na \(R_0\) dává (QA3). Pro bod fyzikálního okolí
\(\rho<R_0/2\) je dostupný polydisk poloměru \(R_0/2\), což dává (QA4).
Žádný nevyčíslený poloměr v těchto místních mezích nezůstává.

Výsledky QA platí pro společný původní příspěvek pólu
a kontinua. Nejde o samostatný odhad každého členu po
jejich rozdělení. Následující dva oddíly doplňují zbývající
regulární místní situace; stacionární vázané a kontinuální
prahy zůstávají odlišnou úlohou.

\section{Původní pulzní matice a přesné analytické rozštěpení}
\hypertarget{analytickuxe9-oznaux10denuxed-obou-pulznuxedch-vux11btvuxed-pux159es-kpi}{%
\subsection{\texorpdfstring{Analytické označení obou pulzních větví přes
\(k=\pi\)}{Analytické označení obou pulzních větví přes k=\textbackslash pi}}\label{analytickuxe9-oznaux10denuxed-obou-pulznuxedch-vux11btvuxed-pux159es-kpi}}

Položme \(k=\pi+h\), kde \(h\) je malá místní souřadnice na kruhu.
Definujme holomorfní funkce s hodnotami určenými v nule:

\[
A(h)=\frac{\sqrt{3+2\cos h}}{\sqrt5},\quad A(0)=1,
\qquad B(h)=-\frac{2\sin(h/2)}{\sqrt5},
\]

\[
\alpha(h)=2\arcsin B(h),\qquad \alpha(0)=0,
\qquad A^2+B^2=1.
\tag{KP13}
\]

Jde o \textbf{podepsanou} analytickou pulzní fázi. Pro reálné \(h\) je
\(\alpha>0\) vlevo od nuly a \(\alpha<0\) vpravo. Oproti původní
nezáporné fázi \(\varepsilon\) se vpravo pouze vymění označení obou
větví. Přitom

\[
\alpha'(h)=-\frac{2\cos(h/2)}{\sqrt{3+2\cos h}},\qquad
\alpha'(0)=-\frac2{\sqrt5}.
\tag{KP14}
\]

Původní pulzní symbol

\[
V(k)=\frac15\begin{pmatrix}
1-4e^{ik}&-2(1+e^{ik})\\
2(1+e^{-ik})&1-4e^{-ik}
\end{pmatrix}
\]

má v těchto souřadnicích vlastní hodnoty \(e^{-i\sigma\alpha(h)}\).
Jejich první kontaktní váhy jsou

\[
H_\sigma(h)=\frac{1-\sigma\alpha'(h)}2,
\qquad H_\sigma(0)=\frac{1+2\sigma/\sqrt5}{2}>0,
\qquad H_++H_-=1.
\tag{KP15}
\]

Váhy jsou holomorfní i v nule. Stejně tak jsou odstranitelné matricové
výrazy pro projektory: čitatel i \(\sin\alpha(h)\) mají společný faktor
\(h\), zatímco \(\alpha'(0)\ne0\). Nevzniká pól projektoru.

Konkrétně pro \(h\ne0\) jsou původní projektory

\[
\mathsf P_\sigma(h)=
\frac{e^{i\sigma\alpha(h)}I-V(\pi+h)}
{2i\sigma\sin\alpha(h)}.
\tag{KP15a}
\]

Čitatel v nule mizí, protože \(V(\pi)=I\). Dělení společným faktorem
\(h\) dává holomorfní matici i v nule. Přesné identity
\(\mathsf P_++\mathsf P_-=I\),
\(V=\sum_\sigma e^{-i\sigma\alpha}\mathsf P_\sigma\),
\(\mathsf P_\sigma\mathsf P_\tau=\delta_{\sigma\tau} \mathsf P_\sigma\)
a \((\mathsf P_\sigma)_{11}=H_\sigma\) se tam prodlouží analyticky. Jde
tedy o skutečné analytické rozštěpení původní pulzní matice.

Aby byla totožnost úplného integrandu doložena již \textbf{před} vnějším
limitem, označme v následující rovnici původní holomorfní objektovou
funkci ve spektrální proměnné \(\lambda\) jako
\(\widehat f_A(P,\lambda)\). Pro reálné \(h,P\) a
\(\lambda=e^{-i(e+i\epsilon)}\), \(\epsilon>0\), je \(|\lambda|>1\).
Holomorfní funkcionální kalkulus původní unitární pulzní matice dává
přesně

\[
\boxed{
\bigl[\widehat f_A(P,\lambda V(\pi+h)^{-1})\bigr]_{11}
=\sum_{\sigma=\pm1}H_\sigma(h)
\widehat f_A(P,\lambda e^{i\sigma\alpha(h)}).
}
\tag{KP15b}
\]

Levá strana je právě původní kontaktní prvek společného objektového a
pulzního operátoru. Všechny skalární argumenty na pravé straně leží mimo
jednotkový kruh; identita proto nepředpokládá žádnou již pokračovanou
hraniční hodnotu ani samostatnou volbu objektové odmocniny.

Úplný původní integrand přesně přechází na

\[
\boxed{
\sum_{\sigma=\pm1}H_\sigma(h)
f_A(K-\pi-h,e-\sigma\alpha(h)).
}
\tag{KP16}
\]

Rovnost platí pro původní vnější hodnoty na obou stranách \(h=0\): vlevo
jsou označení stejná, vpravo jsou obě prohozena. V nule se obě objektové
argumenty shodují a součet vah je 1. Proto lze původní místní periodu
přesně zapsat jako součet integrálů obou analyticky označených větví.
Nedělí se samostatně žádná singularita objektového pólu a jeho kontinua.

\hypertarget{vuxfdmux11bna-pulznuxedch-vux11btvuxed-mimo-objektovou-hranu}{%
\subsection{Výměna pulzních větví mimo objektovou
hranu}\label{vuxfdmux11bna-pulznuxedch-vux11btvuxed-mimo-objektovou-hranu}}

Při \(h=0\) je \(\omega=e\). Podle (PT31) zde v původním okně není
fyzikální objektový pól pro žádné \(K\): v případné mezeře je původní
reálný jmenovatel \(H<0\); na otevřeném kontinuu fyzikální pól rovněž
není. Kompaktní konce jsou regulární při této fázi. Jedinými možnými
místními objektovými hranami jsou \(t=4\), nikoli \(t=-4\).

Pokud \(t\ne4\), jsou oba členy (KP16) na původním listu holomorfní v
\((h,K,e)\) v malém okolí základního bodu. Jejich místní integrál je
tudíž holomorfní v \((K,e)\). Zbývá přesně současná pulzní výměna a
objektová hrana (PT30).

\subsection{Nenulový objektový jmenovatel při současné hraně}
Při \(h=0\), \(v_0=4+5\cos e_0\), \(q_0=1\) a
\(x_0=e^{-i(e_0-g)}\) přímé dosazení do QA5 dává
\[
(D_o)_0=2x_0(x_0-1)
\bigl[\cos(e_0-g)-\cos e_0\bigr].\tag{KP21}
\]
Na celém \(I_e\) platí
\(\cos e_0<-4/5\), \(\cos(e_0-g)>3/10\) a
\(\pi/3<e_0-g<\pi/2\). Proto \(|x_0-1|>1\) a
\[
|(D_o)_0|>11/5.\tag{KP22}
\]
To je středová mez použitá v následujícím číselném důkazu.

\section{Dva číselné obvody při výměně pulzních větví}
\hypertarget{dvux11b-analytickuxe9-pulznuxed-vux11btve}{%
\subsection{Dvě analytické pulzní
větve}\label{dvux11b-analytickuxe9-pulznuxed-vux11btve}}

Pro \(k=\pi+h\) zvolme \[
\alpha(h)=2\arcsin\left(-\frac{2\sin(h/2)}{\sqrt5}\right),\qquad
\alpha(0)=0,
\] \[
\alpha'=-\frac{2\cos(h/2)}{\sqrt{3+2\cos h}},
\qquad H_\sigma(h)=\frac{1-\sigma\alpha'(h)}2.
\tag{QB1}
\]

Jde o podepsanou analytickou fázi. Přesný původní součet obou pulzních
příspěvků je \[
\sum_{\sigma=\pm1}H_\sigma(h)
f_A(K-\pi-h,e-\sigma\alpha(h)).
\tag{QB2}
\] Vlevo od nuly se označení shoduje s původní nezápornou fází, vpravo
jsou obě označení prohozena. Analytické projektory \[
\mathsf P_\sigma(h)=
\frac{e^{i\sigma\alpha(h)}I-V(\pi+h)}
{2i\sigma\sin\alpha(h)}
\] mají odstranitelný společný faktor \(h\). Pro původní vnější
\(\lambda=e^{-i(e+i\epsilon)}\), \(\epsilon>0\), platí přesně \[
[\widehat f_A(P,\lambda V^{-1})]_{11}
=\sum_\sigma H_\sigma(h)
\widehat f_A(P,\lambda e^{i\sigma\alpha(h)}).
\tag{QB3}
\] Proto lze každou ze dvou zachovaných analytických period deformovat
samostatně; nejde o vynechání jednoho pulzního členu.

\hypertarget{stux159edy-a-jednotnuxe9-pulznuxed-meze}{%
\subsection{Středy a jednotné pulzní
meze}\label{stux159edy-a-jednotnuxe9-pulznuxed-meze}}

Fixujme \(e_0\in I_e\) a jeden ze dvou bodů \[
P_0=\pm\arccos(4+5\cos e_0),\qquad K_0=P_0+\pi,\qquad h=0.
\] Na obou větvích je \(\omega_0=e_0\), \(t_0=4\), \(q_0=1\). Původní
jmenovatel má \[
\boxed{|(D_o)_0|>\frac{11}{5}.}
\tag{QB4}
\] Dále \(|v_0|=|4+5\cos e_0|<1/8\): stačí \(v_a<1/16\),
\(|e_0-4\pi/5|\le1/1000\) a \(v_a+1/200<1/8\).

Pro komplexní \(|h|\le1/64\) platí \(|3+2\cos h-5|<1/16\). Jmenovatel v
(QB1) má modul větší než 2. Přímá derivace \[
\alpha''=
\frac{\sin(h/2)}{\sqrt{3+2\cos h}}
-\frac{2\cos(h/2)\sin h}{(3+2\cos h)^{3/2}}
\] dává \[
\boxed{|\alpha'|<2,\quad |\alpha''|<5,\quad
|\alpha(h)|<2|h|,\quad |H_\sigma|<2.}
\tag{QB5}
\] Větve odmocniny a arkussinu jsou určené hodnotami v nule; jejich
argumenty jsou v těchto discích oddělené od příslušných hran.

\hypertarget{tytuxe9ux17e-konkruxe9tnuxed-polomux11bry}{%
\subsection{Tytéž konkrétní
poloměry}\label{tytuxe9ux17e-konkruxe9tnuxed-polomux11bry}}

Zvolme \[
\boxed{r=2^{-24},\qquad R_0=2^{-36},\qquad \rho=2^{-64}.}
\tag{QB6}
\] Pro každé \(\sigma\) položme \[
P=K-\pi-h,\quad \omega_\sigma=e-\sigma\alpha(h),\quad
t_\sigma=\cos P-5\cos\omega_\sigma.
\] V základním bodě je \[
a_\sigma=(t_\sigma)_{h,0}
=\sin P_0+\sigma2\sqrt5\sin e_0,\qquad
1<|a_\sigma|<6,\quad \operatorname{sgn}a_\sigma=\sigma.
\tag{QB7}
\] Dolní mez plyne z \(\sin e_0>1/2\), horní z \(\sqrt5<9/4\).

Na pracovním boxu \(|h|,|K-K_0|,|e-e_0|\le2^{-10}\) jsou podle (QB5) \[
|(t_\sigma)_{hh}|<100,\qquad
|(t_\sigma)_K|<2,\qquad |(t_\sigma)_e|<10.
\] Na kruhu \(|h|=r\), pro parametry v \(R_0\), tedy \[
|t_\sigma-4-a_\sigma h|
\le50r^2+12R_0<r/16.
\tag{QB8}
\] Stejně jako v QA dostáváme na kruhu \[
r/2<|t_\sigma-4|<7r,\quad
7<|t_\sigma+4|<9,\quad
\sqrt r<|\mathcal R_\sigma|<8\sqrt r.
\tag{QB9}
\]

Rouchého věta dává pro každou analytickou pulzní větev právě jeden
jednoduchý bod větvení \(h_\sigma(K,e)\) uvnitř kruhu. Druhá objektová
hrana zde není. Kořen je holomorfní na celém polydisku \(R_0\), reálný
pro reálné parametry a \(|h_\sigma|<r\).

\hypertarget{ux17euxe1dnuxfd-kontaktnuxed-puxf3l-na-kteruxe9mkoli-objektovuxe9m-listu}{%
\subsection{Žádný kontaktní pól na kterémkoli objektovém
listu}\label{ux17euxe1dnuxfd-kontaktnuxed-puxf3l-na-kteruxe9mkoli-objektovuxe9m-listu}}

V celém disku \(|h|\le r\) a pro parametry v \(R_0\) jsou \[
|P-P_0|<2r,\quad |\omega_\sigma-e_0|<3r,\quad
|v-v_0|<4r,\quad |x-x_0|<6r.
\] Horní mez \(|t_\sigma-4|<7r\) platí i uvnitř disku. Proto na obou
objektových listech \[
|q-1|\le\frac{|t_\sigma-4|+|\mathcal R_\sigma|}{4}
<3\sqrt r,\qquad |q|<2.
\tag{QB10}
\] Jsou splněny \(|x|<2,\ |v|<1/4,\ |q|<2\). Původní polynom
\(D_o=x(x^2-\frac25vx-\frac35) +q(\frac35x^2+\frac25vx-1)\) má v tomto
boxu \[
|(D_o)_x|<18,\qquad |(D_o)_v|<4,\qquad |(D_o)_q|<4.
\] Tudíž \[
|D_o-(D_o)_0|<18(6r)+4(4r)+4(3\sqrt r)
=124r+12\sqrt r<13\sqrt r<1.
\] Z (QB4) plyne v celém disku na obou listech \[
\boxed{|D_o|>1.}
\tag{QB11}
\] Původní čitatel tedy splňuje \(|N|<8\) a \[
\boxed{|f_A|<33.}
\tag{QB12}
\] Uvnitř není žádný skrytý pól. Pulzní projektory jsou analytické podle
(QB3), takže nepřidávají další singularitu.

\hypertarget{spruxe1vnuxe9-opaux10dnuxe9-vnux11bjux161uxed-smux11bry}{%
\subsection{Správné opačné vnější
směry}\label{spruxe1vnuxe9-opaux10dnuxe9-vnux11bjux161uxed-smux11bry}}

Ve středu \[
\partial_e h_\sigma
=-\frac{5\sin e_0}{\sin P_0+\sigma2\sqrt5\sin e_0},
\] \[
\operatorname{sgn}\partial_e h_\sigma=-\sigma,\qquad
\frac5{12}<|\partial_e h_\sigma|<5.
\tag{QB13}
\] Z kořenové meze \(|h_\sigma|<r\) na \(R_0\) dávají Cauchyovy odhady
na polovičním polydisku \[
|\partial_Kh_\sigma|,|\partial_eh_\sigma|\le2^{13},
\qquad
|\partial_e\partial_Kh_\sigma|,|\partial_e^2h_\sigma|\le2^{51}.
\] Na \(|K-K_0|,|e-e_0|<\rho\) proto \[
|h_\sigma|\le2^{-50}<r/4,\qquad
|\partial_eh_\sigma-\partial_eh_\sigma(K_0,e_0)|\le2^{-12}.
\tag{QB14}
\] Reálná část energetické derivace si zachovává znaménko \(-\sigma\) a
modul větší než \(1/3\). Při původním \(e+i\epsilon\) tedy \(h_+\) leží
dole a \(h_-\) nahoře.

Perioda větve \(\sigma=+1\) se deformuje na pevný horní půlkruh
\(\Gamma_+\), perioda větve \(\sigma=-1\) na dolní \(\Gamma_-\). Každý
obvod míjí vlastní jediný bod větvení a na obou listech není žádný pól.
Obě deformace zachovávají původní vnější integrály, a proto i jejich
přesný součet (QB3). Opačné posuny kořenů nesmějí být zaměněny za
společné sevření jedné nedělitelné periody.

\hypertarget{konkruxe9tnuxed-uxfaplnuxfd-souux10det-a-derivace}{%
\subsection{Konkrétní úplný součet a
derivace}\label{konkruxe9tnuxed-uxfaplnuxfd-souux10det-a-derivace}}

Položme \[
I_\sigma(K,e)=\int_{\Gamma_\sigma}
H_\sigma(h)f_A(K-\pi-h,e-\sigma\alpha(h))
\frac{\mathrm dh}{2\pi},\qquad
I_{\rm pulse}=I_++I_-.
\] Každý z obou pevných obvodů má správně pokračovaný list holomorfní na
celém komplexním parametrickém polydisku \(R_0\). Z délky \(\pi r\),
(QB5) a (QB12) plyne \[
|I_\sigma|<33r,\qquad
\boxed{|I_{\rm pulse}|<66r<2^{-17}.}
\tag{QB15}
\] Cauchyův vzorec tedy dává ve středu \[
\boxed{
|\partial_K^m\partial_E^n I_{\rm pulse}(K_0,e_0)|
\le m!n!\,2^{36(m+n)-17}.
}
\tag{QB16}
\] Na celém fyzikálně identifikovaném okolí \(\rho\) platí stejná mez s
exponentem \(37(m+n)-17\).

\hypertarget{rozsah}{%
\subsection{Rozsah}\label{rozsah}}

Číselná věta se týká obou průsečíků pulzní výměny s objektovou hranou,
na celém původním uzavřeném \(I_e\). Mimo objektovou hranu je analytické
rozštěpení pulzu přímo regulární, jak dokazuje KP; jednotný globální
odhad zbývající periody tím však ještě není přidělen.

Vnější předpis, původní integrál \(\mathrm dk/(2\pi)\), oba pulzní
členy, objektové kontinuum, původní jev a normalizace jsou zachovány.
Nejde o nový globální čas ani provedený experiment.

\section{Společný číselný obvod u obou kuželových konců}
\hypertarget{vuxfdsledek-se-spoleux10dnuxfdmi-polomux11bry}{%
\subsection{Výsledek se společnými
poloměry}\label{vuxfdsledek-se-spoleux10dnuxfdmi-polomux11bry}}

Fixujme libovolný bod obou kuželových křivek (PT28):

\[
P_0=\pi,\quad \omega_0=\pi,\quad
\varepsilon(k_0)=\pi-e_0,\quad K_0=\pi+k_0,
\qquad e_0\in I_e.
\tag{QC1}
\]

Položme

\[
\boxed{r=2^{-24},\qquad R_0=2^{-36},\qquad \rho=2^{-64}.}
\tag{QC2}
\]

V místní souřadnici \(p=P-\pi\) zvolme půlkruh \(\Gamma:|p|=r\), vedený
od \(-r\) do \(r\) v polorovině opačné ke znaménku
\(u_0=\varepsilon'(k_0)\). Na něm pokračujeme původní vnější objektový
list.

Místní perioda

\[
I_\Gamma(K,e)=\frac1{2\pi}\int_\Gamma
h_-(K-\pi-p)
f_A(\pi+p,e+\varepsilon(K-\pi-p))\,\mathrm dp
\tag{QC3}
\]

je holomorfní na komplexním polydisku \(|K-K_0|,|e-e_0|<R_0\) a splňuje

\[
\boxed{|I_\Gamma|<64r=2^{-18}.}
\tag{QC4}
\]

Pro reálné \(|K-K_0|,|e-e_0|<\rho\) se rovná původnímu vnějšímu
příspěvku intervalu \(-r\le p\le r\), včetně obou stran kuželového bodu
na hybnostním kruhu. Ve středu proto pro \(m,n\ge0\) platí

\[
\boxed{
|\partial_K^m\partial_E^n I_\Gamma(K_0,e_0)|
\le m!n!\,2^{36(m+n)-18}.
}
\tag{QC5}
\]

Na celém fyzikálně identifikovaném okolí poloměru \(\rho\) platí

\[
\boxed{
|\partial_K^m\partial_E^n I_\Gamma(K,e)|
\le m!n!\,2^{37(m+n)-18}.
}
\tag{QC6}
\]

Zde \(\partial_E=\partial_e\). Poloměr \(R_0\) označuje analytické
pokračování na pevném obvodu, zatímco \(\rho\) je doložené okolí shody s
původním fyzikálním předpisem.

\hypertarget{vyux10duxedslenuxe9-holomorfnuxed-pulznuxed-okoluxed}{%
\subsection{Vyčíslené holomorfní pulzní
okolí}\label{vyux10duxedslenuxe9-holomorfnuxed-pulznuxed-okoluxed}}

V kuželových centrech je \(\varepsilon_0=\pi/5+\xi\),
\(|\xi|\le1/1000\). Z \(\sqrt5<56/25\) plyne \(\cos(\pi/5)<81/100\), a
tedy \(|\cos\varepsilon_0|<811/1000\).

Pišme \(a(k)=(1-4\cos k)/5\). Na komplexním disku \(|k-k_0|\le1/64\)
jsou \(|\sin k|,|\cos k|<2\), takže \(|a'|,|a''|<8/5\) a

\[
|a-a(k_0)|\le\frac1{40},\qquad
|a|<\frac{811}{1000}+\frac1{40}
=\frac{209}{250}<\frac67.
\tag{QC7}
\]

Funkce \(b(k)=\sqrt{1-a(k)^2}\) má holomorfní větev s
\(b(k_0)=\sin\varepsilon_0>0\). Její modul je větší než
\(\sqrt{13}/7>1/2\). Původní fáze se proto holomorfně pokračuje a

\[
\varepsilon'=-a'/b,\qquad
\varepsilon''=-a''/b-a(a')^2/b^3.
\]

Z těchto přesných identit stačí konzervativní komplexní meze

\[
\boxed{|\varepsilon'|<4,\qquad |\varepsilon''|<24.}
\tag{QC8}
\]

První pravá strana se omezuje již \(16/5<4\); druhá
\(16/5+(6/7)(64/25)8<24\). Původní reálná středová rychlost současně
splňuje mnohem ostřejší doložené meze

\[
4/5<|u_0|\le2/\sqrt5.
\tag{QC9}
\]

\hypertarget{dvux11b-analytickuxe9-hranovuxe9-funkce-a-jejich-oddux11blenuxed-na-obvodu}{%
\subsection{Dvě analytické hranové funkce a jejich oddělení na
obvodu}\label{dvux11b-analytickuxe9-hranovuxe9-funkce-a-jejich-oddux11blenuxed-na-obvodu}}

Položme

\[
\nu(p,K,e)=e+\varepsilon(K-\pi-p)-\pi,
\qquad
f_j(p,K,e)=\sin(p/2)-j\sqrt5\sin(\nu/2),\quad j=\pm1.
\]

Původní práh se přesně rozkládá jako

\[
t-4=2f_+f_-,\qquad t=-\cos p+5\cos\nu.
\tag{QC10}
\]

Na pracovním boxu \(|p|,|K-K_0|,|e-e_0|\le2^{-10}\) platí
\(|k-k_0|\le2^{-9}\) a \(|\nu|\le9\cdot2^{-10}<1/64\). Přímé derivace za
použití (QC8) a
\(|\sin(p/2)|,|\cos(p/2)|,|\sin(\nu/2)|,|\cos(\nu/2)|<2\) dávají

\[
|\partial_p^2 f_j|<100,\qquad
|\partial_Kf_j|<10,\qquad |\partial_ef_j|<3.
\tag{QC11}
\]

Například první mez je menší než
\(1/2+\sqrt5[2(4^2)/4+2(24)/2]<81<100\).

Středový koeficient je

\[
a_j=\partial_pf_j(0,K_0,e_0)=\frac{1+j\sqrt5u_0}{2},
\qquad \frac38<|a_j|\le\frac32.
\tag{QC12}
\]

Pro dolní mez stačí (QC9) a \(\sqrt5>35/16\). Na celém kruhu \(|p|=r\) a
parametrickém polydisku \(R_0\) tedy

\[
|f_j-a_jp|\le50r^2+13R_0<r/16.
\tag{QC13}
\]

Poslední porovnání je přesně \(50/2^{24}+13/2^{12}<1/16\). Z toho plyne

\[
\boxed{5r/16<|f_j|<2r\qquad(|p|=r).}
\tag{QC14}
\]

Rouchého věta dává pro každý \(f_j\) právě jeden jednoduchý kořen
\(p_j(K,e)\) uvnitř kruhu. Každý je holomorfní na celém parametrickém
polydisku \(R_0\); lze jej vyjádřit konturovým integrálem
\(p\,\partial_pf_j/f_j\). Pro reálné parametry je reálný, protože
příslušná \(f_j\) je reálná analytická funkce a kořen v symetrickém
disku je jediný. Kořeny obou různých faktorů se mohou vzájemně setkat,
aniž by kterýkoli přestal být jednoduchým kořenem svého vlastního
faktoru.

Na obvodu proto

\[
r^2/8<|t-4|<8r^2,
\qquad 7<|t+4|<9,
\qquad r/2<|\mathcal R|<9r,
\quad \mathcal R^2=t^2-16.
\tag{QC15}
\]

Princip maxima dává horní meze \(|f_j|<2r\), \(|t-4|<8r^2\) i uvnitř
disku. Tam nejsou jiné objektové odmocninové nuly: faktor \(t+4\) je
nenulový a každý \(f_j\) má právě svůj jediný kořen.

\hypertarget{pux16fvodnuxed-jmenovatel-je-nenulovuxfd-na-obou-listech-v-celuxe9m-disku}{%
\subsection{Původní jmenovatel je nenulový na obou listech v celém
disku}\label{pux16fvodnuxed-jmenovatel-je-nenulovuxfd-na-obou-listech-v-celuxe9m-disku}}

Používáme původní polynomy (QA5). Ve středu je \(v_0=-1\), \(q_0=1\),
\(x_0=-e^{ig}\) a

\[
(D_o)_0=(x_0+1)^2(x_0-1),\qquad
|(D_o)_0|=\sqrt5>11/5.
\tag{QC16}
\]

Na celém disku \(|p|\le r\) a polydisku \(R_0\) platí

\[
|\nu|\le4r+5R_0<5r,\quad
|x-x_0|<10r,\quad |v+1|<r^2,
\]

\[
|q-1|=\frac14|t-4-\mathcal R|<3r.
\tag{QC17}
\]

Poslední mez používá \(|\mathcal R|<9r\), které je platné na
\textbf{obou} listech i uvnitř disku. Zvláště jsou \(|x|,|v|,|q|<2\).
Přímé derivace původního polynomu

\[
D_o=x^3-\left(\frac25v-\frac35q\right)x^2
+\left(-\frac35+\frac25vq\right)x-q
\]

dávají na tomto polydisku nezávislých proměnných

\[
|(D_o)_x|<24,\qquad |(D_o)_v|<4,\qquad |(D_o)_q|<5.
\]

Tudíž

\[
|D_o-(D_o)_0|<24(10r)+4r^2+5(3r)<256r<1/5,
\]

\[
\boxed{|D_o|>2.}
\tag{QC18}
\]

V celém uvažovaném algebraickém okolí není žádný kontaktní pól. Nejde
pouze o mez na samotném obvodu. Pulzní fáze a její projektor jsou
holomorfní podle (QC7)--(QC8).

\hypertarget{vyux10duxedslenuxfd-spoleux10dnuxfd-fyzikuxe1lnuxed-smux11br}{%
\subsection{Vyčíslený společný fyzikální
směr}\label{vyux10duxedslenuxfd-spoleux10dnuxfd-fyzikuxe1lnuxed-smux11br}}

Z holomorfních kořenů z (QC13)--(QC14) máme na polydisku \(R_0\) \(|p_j|<r\).
Ve středu jejich přesné derivace jsou

\[
\partial_ep_j(K_0,e_0)=\frac1{u_0+j/\sqrt5},\qquad
\frac23<|\partial_ep_j|<3,
\qquad \operatorname{sgn}(\partial_ep_j)=\operatorname{sgn}u_0.
\tag{QC19}
\]

Cauchyovy meze na polovičním polydisku jsou

\[
|\partial_Kp_j|,|\partial_ep_j|\le2r/R_0=2^{13},
\]

\[
|\partial_e\partial_Kp_j|,|\partial_e^2p_j|
\le8r/R_0^2=2^{51}.
\tag{QC20}
\]

Na malém komplexním polydisku \(\rho=2^{-64}\) tedy

\[
|p_j|\le2^{14}\rho=2^{-50}<r/4,
\]

\[
\left|\partial_ep_j-\frac1{u_0+j/\sqrt5}\right|
\le2^{52}\rho=2^{-12}.
\tag{QC21}
\]

Reálné části obou energetických derivací mají stejné znaménko jako
\(u_0\) a absolutní hodnotu větší než \(1/2\). Pro reálné parametry v
malém okolí a kladné \(\epsilon\), které ponechá \((K,e+i\epsilon)\)
uvnitř téhož polydisku, integrace po svislé úsečce dává

\[
\boxed{
\operatorname{sgn}\operatorname{Im}p_+(K,e+i\epsilon)
=\operatorname{sgn}\operatorname{Im}p_-(K,e+i\epsilon)
=\operatorname{sgn}u_0.
}
\tag{QC22}
\]

V disku jsou pouze tyto dva objektové body větvení, případně jejich
společný dvojnásobný průmět, a žádný pól. Reálný interval se proto může
deformovat na pevný půlkruh \(\Gamma\) v opačné polorovině, bez
překročení singularity. Po celou deformaci se zachovává původní vnější
objektový list. Limit \(\epsilon\downarrow0\) dává přesně původní úplný
kuželový příspěvek.

\hypertarget{amplituda-periody-a-vux161echny-koneux10dnuxe9-derivace}{%
\subsection{Amplituda periody a všechny konečné
derivace}\label{amplituda-periody-a-vux161echny-koneux10dnuxe9-derivace}}

Z \(|x|,|v|,|q|<2\) plyne

\[
|N|<4+\frac35(2+2)2+4<13.
\]

Podle (QC8),(QC18) máme \(|h_-|<3\) a \(|f_A|<27\). Na obvodu žádný
radikál nenuluje a jmenovatel je nenulový. Původní list lze po oblouku a
v jednoduše souvislém parametrickém polydisku \(R_0\) pokračovat
společně. Integrand a jeho perioda jsou tam holomorfní. Z délky
\(\pi r\) plyne

\[
|I_\Gamma|
\le\frac{\pi r}{2\pi}\,3\cdot27
=\frac{81}{2}r<64r=2^{-18}.
\tag{QC23}
\]

Cauchyův vzorec na polydisku \(R_0\) dává (QC5). V bodě fyzikálního
okolí \(\rho<R_0/2\) je k dispozici polydisk poloměru \(R_0/2\), což
dává (QC6).

Tím jsou pro obě kuželové křivky vyčísleny obvod, správné jednostranné
pokračování, místní amplituda a libovolný konečný počet parametrických
derivací. Ostatní části úplného \(F_2\), stacionární prahy a původní
ostré energetické konce nejsou těmito místními konstantami nahrazeny.
Původní pravděpodobnostní jev, normalizace, všech 26 štítků a
paměť/reference se nemění.

\section{Vyčíslené komplexní okolí celého prvního kontaktu}
\hypertarget{pux16fvodnuxed-pux159esnuxfd-integruxe1l}{%
\subsection{Původní přesný
integrál}\label{pux16fvodnuxed-pux159esnuxfd-integruxe1l}}

Pišme \(e=E-\beta\), \(g=2\pi/5\), \(e_0\in I_e\). Původní vzorec je \[
F_1(E)=\int_{-\pi}^{\pi} f_B(P,e)\frac{\mathrm dP}{2\pi},
\quad
f_B(P,e)=\frac{\mathcal R-4i\sin(e-g)}
{t+4\cos(e-g)},
\quad
t=\cos P-5\cos e,\quad \mathcal R^2=t^2-16.
\tag{FQ1}
\] Odmocnina se určuje původním vnějším předpisem \(e+i0\). Na reálném
původním okně platí dříve doložené meze \[
3<t<51/10,\qquad d_B=t+4\cos(e-g)>21/5.
\] Proto \(|\mathcal R|<16/5\), \(|f_B|<12/7<2\) a \(|F_1|<2\). Na
\(|P|\le1\) je \(t>9/2\), \(\mathcal R>2\), tedy \(\Re f_B>20/91\). Z
kladnosti celé kontaktní spektrální míry plyne \[
\boxed{\Re F_1(E)>20/(91\pi)>1/15,\qquad |F_1(E)|<2.}
\tag{FQ2}
\]

\hypertarget{jeden-vuxfdslovnuxfd-obvod-pro-celuxfd-hybnostnuxed-kruh}{%
\subsection{Jeden výslovný obvod pro celý hybnostní
kruh}\label{jeden-vuxfdslovnuxfd-obvod-pro-celuxfd-hybnostnuxed-kruh}}

Zvolme \[
a=2^{-8},\qquad b=2^{-16},\qquad
P(\theta)=\theta-ia\sin\theta,\quad -\pi\le\theta\le\pi.
\tag{FQ3}
\] Pro \(|e-e_0|<b\) jsou všechny níže uvedené goniometrické funkce
holomorfní a \(|\sin z|,|\cos z|<2\) na použitých krátkých úsečkách. Ve
středu je \(|4+5\cos e_0|<1/8\), jak bylo doloženo pro QB.

Je-li \(|\cos\theta|\le1/2\), pak přesně \[
\operatorname{Im}\cos P(\theta)
=\sin\theta\,\sinh(a\sin\theta)\ge a\sin^2\theta\ge3a/4.
\] Proto \[
\operatorname{Im}t\ge3a/4-10b>a/2.
\] Je-li \(|\cos\theta|>1/2\), pak \[
|\cos\theta-(4+5\cos e_0)|>3/8.
\] Změna způsobená deformací obvodu a energií je nejvýše
\(2a+10b<1/16\), takže také zde \(|t-4|>a/2\). Celkem na celém obvodu \[
\boxed{|t-4|>2^{-9},\qquad |t+4|>6,\qquad |t|<7.}
\tag{FQ4}
\] Pro druhou mez stačí původní \(t_0>3\) a změna menší než \(1/16\).
Odtud na obou listech \(|\mathcal R|<9\).

Původní jmenovatel má reálnou středovou mez \(d_{B,0}>21/5\). Jeho změna
je menší než \(2a+18b<1/5\). Tudíž \[
\boxed{|d_B|>4,\qquad |f_B|<17/4<5.}
\tag{FQ5}
\] Na obvodu nejsou žádné objektové body větvení ani kontaktní póly.

\hypertarget{shoda-s-pux16fvodnuxedm-vnux11bjux161uxedm-pux159edpisem}{%
\subsection{Shoda s původním vnějším
předpisem}\label{shoda-s-pux16fvodnuxedm-vnux11bjux161uxedm-pux159edpisem}}

Použijme homotopii \(P_s(\theta)=\theta-isa\sin\theta,\ 0\le s\le1\).
Pro \(e=e_r+i\epsilon\), kde \(\epsilon>0\), \(e_r\) je reálné a
\(|e-e_0|<b\), je \[
\operatorname{Im}t(P_s,e)
=\sin\theta\,\sinh(sa\sin\theta)
+5\sin e_r\,\sinh\epsilon>0.
\tag{FQ6}
\] Používá se \(\sin e_r>0\), zaručené celým daným oknem. Radikál tedy
během deformace nikdy nenuluje a (FQ5) ve slabé podobě pro \(sa\)
vylučuje také kontaktní pól. Původní vnější objektový list lze
pokračovat po celé homotopii a integrál se nemění.

Na pevném obvodu \(P(\theta)\) se tento list holomorfně pokračuje na
celý parametrický disk \(|e-e_0|<b\). Jeho periodicita v \(\theta\) se
zachovává z původního listu; žádná nula radikálu nemůže změnit jeho
monodromii. Perioda tím dává analytické pokračování celého původního
\(F_1\).

Z \(|P'(\theta)|\le1+a<5/4\) a (FQ5) dostáváme \[
\boxed{|F_1(e)|<8\qquad(|e-e_0|<2^{-16}).}
\tag{FQ7}
\]

\hypertarget{vyux10duxedslenuxe1-komplexnuxed-kladnost-a-derivace}{%
\subsection{Vyčíslená komplexní kladnost a
derivace}\label{vyux10duxedslenuxe1-komplexnuxed-kladnost-a-derivace}}

Na polovičním energetickém disku Cauchyův vzorec dává \[
|F_1'(e)|\le 2\cdot8/b=2^{20}.
\] Zvolme menší společný komplexní poloměr \[
\boxed{\rho_1=2^{-30}.}
\tag{FQ8}
\] Pro \(|e-e_0|<\rho_1\) je \[
|F_1(e)-F_1(e_0)|<2^{20}\rho_1=2^{-10}.
\] Protože \(1/15-1/16=1/240>2^{-10}\), z (FQ2) plyne \[
\boxed{\Re F_1(e)>1/16,\qquad |F_1(e)|<3
\quad\text{pro }|e-e_0|<2^{-30}.}
\tag{FQ9}
\]

Ve středu jsou pro \(j\ge0\) doložené derivace \[
\boxed{|F_1^{(j)}(e_0)|\le j!\,2^{16j+3}.}
\tag{FQ10}
\] Na celém menším okolí \(\rho_1<b/2\) platí \[
|F_1^{(j)}(e)|\le j!\,2^{17j+3}.
\] Jde o celý první kontakt. K přesné kvantifikaci úplného přenosového
sloupce u vázaného prahu tím přestává být neznámou šířka analytického
okolí \(F_1\). Zbytek periody druhého kontaktu \(F_2\) zůstává
samostatnou úlohou.

\section{Přesný kvantitativní tvar obou vázaných prahů}
\hypertarget{vuxfdsledek-a-pux159esnuxe9-vymezenuxed-jeho-dosahu}{%
\subsection{Výsledek a přesné vymezení jeho
dosahu}\label{vuxfdsledek-a-pux159esnuxe9-vymezenuxed-jeho-dosahu}}

Nechť \((P_0,K_0,E_0)\) je libovolný z obou fyzikálních stacionárních
bodů (BS25), s \(E_0\) v původním okně
\(I=[8\pi/5-10^{-3},8\pi/5+10^{-3}]\). Všechny následující konstanty
jsou společné pro oba body a celé toto okno. Pro příslušnou nižší
objektovou větev a pulzní znaménko \(\sigma=-1\) pišme

\[
\Phi(P,K)=\beta+\omega_b(P)-\varepsilon(K-P),\qquad
\beta=4\pi/5.
\tag{BF1}
\]

Existují holomorfní funkce \(P_b(K)\), \(E_b(K)=\Phi(P_b(K),K)\) a
přesně vymezený místní příspěvek původního integrálu takové, že pro

\[
|K-K_0|<2^{-88},\qquad
Z=E-E_b(K),\qquad |Z|<2^{-152}
\tag{BF2}
\]

platí na vnější fyzikální větvi

\[
\boxed{
F_{2,\mathrm{loc}}(K,E)
=\frac{A(K,Z)}{\sqrt Z}+B_{\mathrm{loc}}(K,Z),
\qquad
2^{-8}<|A(K,Z)|<2^{12},\quad
|B_{\mathrm{loc}}(K,Z)|<2^{80}.}
\tag{BF3}
\]

Obě funkce \(A,B_{\mathrm{loc}}\) jsou holomorfní v obyčejných
proměnných \((K,Z)\), včetně \(Z=0\). Vnější odmocnina je určena pro
\(\operatorname{Im}Z>0\) kladnou reálnou i imaginární částí. Její
pokračování přes zápornou reálnou osu odpovídá \(0<\arg Z<2\pi\).
Zejména

\[
\sqrt{Z+i0}=\sqrt Z\quad(Z>0),\qquad
\sqrt{Z+i0}=i\sqrt{|Z|}\quad(Z<0).
\tag{BF4}
\]

Nenulový hlavní koeficient má původní normalizaci

\[
\boxed{
A(K_0,0)=h_{-}(K_0-P_0)\,r_b(P_0)
\sqrt{\frac{2}{\Phi_{PP}(P_0,K_0)}}>\frac3{256}.}
\tag{BF5}
\]

Obecně není dovoleno nahradit \(A(K,Z)\) konstantou \(A(K,0)\) a zbytek
nazvat analytickým: zůstaly by členy \(Z^{1/2},Z^{3/2},\ldots\). Tvar
(BF3) již obsahuje všechny tyto půlcelé mocniny. Logaritmický člen u
tohoto jednoduchého jednorozměrného kvadratického záhybu nevzniká.

Zároveň se pro \textbf{úplný} původní přenosový sloupec, na reálné
spektrální straně \(0<Z<2^{-152}\) a při \(E\in I\), dokáže bez odhadu
komplementu následující mez pro jeho okrajové hodnoty (téměř všude,
popřípadě s příslušným mezním významem):

\[
\boxed{
\left\|
\frac{(-2F_2,4)^T}{4+F_1F_2}
-\binom{-2/F_1}{0}
\right\|_2\le 2^{20}\sqrt Z.}
\tag{BF6}
\]

Pro oboustranný holomorfní odmocninový tvar celého sloupce je navíc
potřeba holomorfie zbývající části fyzikální periody a její číselná mez.
Potřebné komplexní okolí úplného \(F_1\) je již samostatně doloženo
(FQ1)--(FQ10). Zbývající požadavek na komplement je u rovnic (BF33)--(BF37) výslovně
oddělen od již dokázaného (BF3),(BF6).

\hypertarget{spoleux10dnuxe1-fyzikuxe1lnuxed-mezera-a-derivace-objektovuxe9ho-puxf3lu}{%
\subsection{Společná fyzikální mezera a derivace objektového
pólu}\label{spoleux10dnuxe1-fyzikuxe1lnuxed-mezera-a-derivace-objektovuxe9ho-puxf3lu}}

V tomto oddílu značí \(X=e^{-i(\omega-g)}\), \(g=2\pi/5\), \(v=\cos P\);
pozdější Morseova souřadnice bude označena \(x\). Úplná fyzikální funkce
objektu je

\[
t=v-5\cos\omega,\qquad R^2=t^2-16,\qquad q=(t-R)/4,
\] \[
N=X^2+\frac35(v+q)X+vq,
\] \[
D=X^3-\left(\frac25v-\frac35q\right)X^2
+\left(-\frac35+\frac25vq\right)X-q,
\qquad f_A=\frac{2XN-D}{D}.
\tag{BF7}
\]

Přesný fyzikální předpis pro R je ten z původního vnějšího resolventu. V
reálné mezeře je \(R=i\operatorname{sgn}(\sin\omega)\sqrt{16-t^2}\). V
každém našem středu platí \(|t_0|<2\) podle (BS28).

Nejprve doplníme horní mez společné křivosti. Z (BS18),(BS19) a
\(v=\cos\delta-4z\sin\delta\) plyne

\[
|v''(\delta)|<1+\frac53+\frac{32}3+\frac{28}3=\frac{68}3,
\quad |v'|>4,
\] \[
|\omega_{vv}|<\frac{17}{48},\qquad
|\omega_A''(p)|<\frac14+\frac{17}{48}=\frac{29}{48}.
\]

Protože na reálném pulzním pásu \(|\varepsilon''|\le1\), společně s
(BS22) dostáváme

\[
\boxed{\frac14<a_0:=\Phi_{PP}(P_0,K_0)<2.}
\tag{BF8}
\]

Pro dolní mez jednoduchosti použijme reálné H z původního jmenovatele.
Na horním reprezentantu pole jsou všechny čtyři členy v \(H_\omega\)
záporné a mezi nimi je \(-4\cos\delta\cos\theta\). Z
\(\cos\delta>17/18\), \(\cos\theta>2/3\) proto plyne

\[
|H_\omega|>\frac{68}{27},\qquad
|D_\omega|=\frac45|H_\omega|>\frac{272}{135}>2.
\tag{BF9}
\]

Rovnost pro D se bere v jeho nule, takže derivace jeho nenulového
fázového faktoru další člen nepřidává. Poloperiodová fyzikální symetrie
(BS1) zachovává absolutní hodnotu této derivace i reziduum; (BF9) tudíž
platí pro oba skutečné nižší stacionární póly.

Vezměme komplexní polydisk

\[
|P-P_0|<R_0=2^{-6},\qquad |\omega-\omega_0|<R_0,
\qquad \omega_0=\omega_b(P_0).
\tag{BF10}
\]

Na něm \(|v|,|X|<2\) a elementární exponenciální odhady dávají

\[
|t-t_0|\le2|P-P_0|+10|\omega-\omega_0|<3/16,
\quad |t|<35/16,
\] \[
|(16-t^2)-(16-t_0^2)|<\frac{67}{16}\frac3{16}
=\frac{201}{256}<1.
\tag{BF11}
\]

Veličina \(16-t_0^2\) je reálná a větší než 12. Její fyzikální odmocnina
má proto na celém (BF10) jedinou zvolenou holomorfní větev s
\(3<|\sqrt{16-t^2}|<5\). Tak \(|q|<2\), a přímo z (BF7)

\[
|N|<16,\qquad |D|<32,\qquad |2XN-D|<128.
\tag{BF12}
\]

Nejde tedy o abstraktní předpoklad komplexní mezery: její konkrétní
šířku zde zajišťuje (BF10),(BF11) pro původní fyzikální odmocninu.

Cauchyho odhady na polovičním polydisku (BF10), s indexy
\(i,j,k\in\{P,\omega\}\), dávají

\[
|D_i|\le2^{12},\qquad |D_{ij}|\le2^{20},\qquad
|D_{ijk}|<2^{29}.
\tag{BF13}
\]

Faktory \(2!\), \(3!\) jsou v těchto číslech zahrnuty; například
\(32\,3!\,128^3<2^{29}\).

\hypertarget{prvnuxed-rouchuxe9ho-krok-celuxfd-lokuxe1lnuxed-laurentux16fv-rozklad}{%
\subsection{První Rouchého krok: celý lokální Laurentův
rozklad}\label{prvnuxed-rouchuxe9ho-krok-celuxfd-lokuxe1lnuxed-laurentux16fv-rozklad}}

Normalizujme \(\mathcal D=D/D_\omega(P_0,\omega_0)\). Ve středu
\(\mathcal D_\omega=1\) a \(|\mathcal D_P|=|\omega_b'(P_0)|<1/4\) podle
(BS16),(BS25). Její druhé derivace jsou omezeny \(2^{19}\), třetí
\(2^{28}\). Použijme nový polydisk s poloměrem

\[
R_o=2^{-24}
\tag{BF14}
\]

v obou proměnných. V něm je změna kterékoli první derivace
\(\mathcal D\) nejvýše \(2^{19}(2R_o)=1/16\). Na kružnici
\(|\omega-\omega_0|=R_o\), pro \(|P-P_0|<R_o\), proto

\[
|\mathcal D-(\omega-\omega_0)|
\le \frac14R_o+2^{20}R_o^2
=\frac5{16}R_o<R_o.
\tag{BF15}
\]

Rouchého věta dává právě jednu nulu v tomto kruhu; derivace
\(|\mathcal D_\omega|>15/16\) je nenulová. Jediná nula tedy tvoří
holomorfní \(\omega_b(P)\). Implicitní diferenciace dává

\[
|\omega_b'|<\frac{1/4+1/16}{1-1/16}=\frac13,
\qquad |\omega_b(P)-\omega_0|\le\frac13|P-P_0|,
\] \[
|\omega_b''|<2^{20},\qquad |\omega_b'''|<2^{43}.
\tag{BF16}
\]

Pro poslední dvě nerovnosti lze v implicitních vzorcích použít

\[
\frac{2^{19}(4/3)^2}{15/16}<2^{20},
\] \[
\frac{2^{28}(4/3)^3+
3\,2^{19}(4/3)2^{20}}{15/16}<2^{43}.
\]

Faktorizace \(D=(\omega-\omega_b(P))d(P,\omega)\) má \(|d|>1\):
integrální vzorec pro d je průměrem \(D_\omega\) na úsečce od
\(\omega_b(P)\) k \(\omega\), přičemž
\(D_\omega/D_\omega(P_0,\omega_0)\) leží v disku \(|z-1|\le1/16\).
Funkce \(J=(2XN-D)/d\) je tedy holomorfní a \(|J|<128\).

Definujme \(2ir_b(P)=J(P,\omega_b(P))\). Je to právě původní spektrální
reziduum, nikoli nová normalizace. Dostáváme přesně

\[
\boxed{
f_A(P,\omega)=\frac{2ir_b(P)}{\omega-\omega_b(P)}
+R_A(P,\omega),\qquad |r_b|<64.}
\tag{BF17}
\]

Na menším oboru \(|P-P_0|<R_o/2\), \(|\omega-\omega_b(P)|<R_o/4\) je
\(R_A\) holomorfní a

\[
|R_A|<2^{34}.
\tag{BF18}
\]

Skutečně \(R_A\) je dělený rozdíl J v proměnné \(\omega\). Z (BF16) leží
příslušné úsečky uvnitř \(|\omega-\omega_0|<5R_o/12\). Cauchyho kruh o
poloměru \(R_o/2\) zůstává v původním oboru J a dává dokonce
\(|J_\omega|\le256/R_o=2^{32}\). Rezerva v (BF18) bude použita beze
změny níže.

\hypertarget{druhuxfd-rouchuxe9ho-krok-stacionuxe1rnuxed-graf-a-jeho-energie}{%
\subsection{Druhý Rouchého krok: stacionární graf a jeho
energie}\label{druhuxfd-rouchuxe9ho-krok-stacionuxe1rnuxed-graf-a-jeho-energie}}

Ze zpřesnění v (CS4) je \(|k_0|<4/15\), kde \(k_0=K_0-P_0\). Na
komplexním disku \(|k-k_0|<1/64\) má \(W(k)=3-2\cos k\) vzdálenost od
své kladné reálné středové hodnoty menší než \(1/16\); zejména
\(|W|>15/16\). Původní pulzní větev má holomorfní fázi s derivacemi

\[
\varepsilon'(k)=\frac{-2\sin(k/2)}{\sqrt{3-2\cos k}},\qquad
\varepsilon''(k)=\frac{-\cos(k/2)}{(3-2\cos k)^{3/2}}.
\]

Protože \(|k|<1/3\), lze použít \(|\sin(k/2)|<1/4\), \(|\cos(k/2)|<2\).
Dostáváme bezpečně \(|\varepsilon'|<2\), \(|\varepsilon''|<8\); na
polovičním disku pak Cauchyho odhad dává \(|\varepsilon'''|<2^{12}\). S
(BF16) z toho plyne, že všechny třetí parciální derivace \(\Phi(P,K)\)
na zde používaných menších polydiscích jsou omezeny \(2^{44}\).

Ve středu má funkce \(\mathcal S=\Phi_P/a_0\) hodnotu nula,
\(\mathcal S_P=1\), \(|\mathcal S_K|=|\varepsilon''(k_0)|/a_0<4\), a
všechny její druhé derivace jsou menší než \(2^{46}\). Zvolme

\[
r_P=2^{-54},\qquad r_K=2^{-58}=r_P/16.
\tag{BF19}
\]

Na P-kružnici o poloměru r\_P je lineární K-člen menší než \(r_P/4\),
zatímco druhý řád je menší než

\[
\tfrac12 2^{46}(r_P+r_K)^2<r_P/128.
\]

Rouchého věta tak opět dává jediný holomorfní kořen \(P_b(K)\). Změny
prvních derivací \(\mathcal S\) jsou menší než
\(2^{46}(r_P+r_K)<1/128\), odkud

\[
|P_b'(K)|<8,\qquad
|P_b(K)-P_0|\le8|K-K_0|,
\qquad |E_b'(K)|<8.
\tag{BF20}
\]

V poslední nerovnosti lze přímo použít
\(E_b'=\Phi_K=-\varepsilon'(K-P_b)\); zápis s 8 ponechává společnou
jednoduchou rezervu. Pro reálné K jsou \(P_b,E_b\) reálné, což plyne z
jednoznačnosti a reálnosti původních fází.

\hypertarget{kvantitativnuxed-holomorfnuxed-morseova-souux159adnice}{%
\subsection{Kvantitativní holomorfní Morseova
souřadnice}\label{kvantitativnuxed-holomorfnuxed-morseova-souux159adnice}}

Omezme se nejprve na \(|K-K_0|<2^{-60}\), \(|P-P_b(K)|<r_M=2^{-58}\).
Definujme

\[
C_M(P,K)=\int_0^1(1-u)
\Phi_{PP}\bigl(P_b(K)+u(P-P_b(K)),K\bigr)\,du.
\]

Taylorova formule a \(\Phi_P(P_b,K)=0\) dávají přesně

\[
\Phi(P,K)-E_b(K)=(P-P_b(K))^2 C_M(P,K).
\tag{BF21}
\]

Z (BF20) a meze \(2^{44}\) pro třetí derivace dostáváme

\[
|C_M-a_0/2|
\le\tfrac12 2^{44}(8+4+1)2^{-60}
=13\,2^{-17}<2^{-13}.
\]

Tedy

\[
1/16<|C_M|<2,\qquad |(C_M)_P|<2^{42}.
\tag{BF22}
\]

Pro druhý odhad je integrál \(\int_0^1u(1-u)du=1/6\). Hodnota C\_M leží
v pevném disku kolem kladného \(a_0/2\), který neobsahuje nulu. Zvolme
tedy holomorfní odmocninu kladnou pro reálné argumenty a položme

\[
x=(P-P_b(K))\sqrt{C_M(P,K)}.
\tag{BF23}
\]

Platí \(|\sqrt{C_M}|>1/4\) a \(|\partial_P\sqrt{C_M}|<2^{43}\). Na
kružnici \(|P-P_b|=r_M\) se (BF23) liší od svého nenulového lineárního
členu o nejvýše \(r_M2^{-15}\). Pro \(|x|<r_M/8=2^{-61}\) je lineární
Rouchého mez nejméně \(r_M/8\). Získáváme jedinou holomorfní inverzi
P=P(x,K) na

\[
|x|<2^{-61},\qquad |K-K_0|<2^{-60},
\] \[
\boxed{
\Phi(P(x,K),K)=E_b(K)+x^2,\quad
|P-P_b|\le4|x|,\quad |P_x|<8.}
\tag{BF24}
\]

Poslední nerovnost plyne z \(|x_P|\ge1/4-r_M2^{43}>1/8\). Na reálném
řezu je inverze orientována kladně; zejména \(P_x(0,K_0)=\sqrt{2/a_0}\).

Tento krok používá skutečnou nenulovou \textbf{společnou} křivost
\(\Phi_{PP}\). Samotná jednoduchost objektového pólu by (BF21)--(BF24)
neimplikovala.

\hypertarget{pux159esnuxfd-integruxe1l-na-fyzikuxe1lnuxed-vux11btvi}{%
\subsection{Přesný integrál na fyzikální
větvi}\label{pux159esnuxfd-integruxe1l-na-fyzikuxe1lnuxed-vux11btvi}}

Zvolme pevný Morseův poloměr

\[
r=2^{-64},\qquad -r\le x\le r.
\tag{BF25}
\]

Místní příspěvek je původní člen s \(\sigma=-1\) integrovaný na reálném
intervalu \(P(-r,K)\le P\le P(r,K)\) s původní mírou
\(\mathrm dP/(2\pi)\). Pro komplexní K se tentýž interval holomorfně
deformuje podle P(x,K). Nejprve je toto určení shodné s původním
integrálem pro reálné K a \(\operatorname{Im}E>0\); až poté se použije
analytické pokračování. Nevkládá se svévolný obchvat pólu ani se
nepřičítá reziduum.

Na tomto intervalu je \(\omega-\omega_b(P)=E-\Phi(P,K)=Z-x^2\). Pro
(BF2) je tedy podmínka (BF18) splněna s velkou rezervou. Označme

\[
a(x,K)=h_-(K-P(x,K))\,r_b(P(x,K))\,P_x(x,K).
\tag{BF26}
\]

Dosazení (BF17) do původní míry \(\mathrm dP/(2\pi)\) dává přesně

\[
F_{2,\mathrm{loc}}
=\frac{i}{\pi}\int_{-r}^r\frac{a(x,K)}{Z-x^2}\,dx
+\int_{-r}^r
h_-R_A\,P_x\frac{dx}{2\pi}.
\tag{BF27}
\]

Faktor \(i/\pi\) je tedy výsledkem původního rezidua \(2ir_b\), nikoli
dodatečné spektrální renormalizace.

Lichá část a má v prvním integrálu nulový příspěvek. Sudou část pišme
\(a_{\rm ev}(x,K)=\sum_{j\ge0}a_{2j}(K)x^{2j}\) a definujme

\[
A(K,Z)=\sum_{j\ge0}a_{2j}(K)Z^j
=\frac{a(\sqrt Z,K)+a(-\sqrt Z,K)}2.
\]

Součet je holomorfní v Z i v nule a nezávisí na znaménku pomocné
odmocniny. Dělený rozdíl

\[
H_*(x^2,Z,K)=\frac{a_{\rm ev}(x,K)-A(K,Z)}{x^2-Z}
\]

je rovněž holomorfní včetně diagonály \(x^2=Z\). Přesný skalární
integrál je

\[
\int_{-r}^r\frac{dx}{Z-x^2}
=-\frac{i\pi}{\sqrt Z}+J_r(Z),\qquad
J_r(Z)=\frac2r\sum_{j=0}^{\infty}
\frac{(Z/r^2)^j}{2j+1}.
\tag{BF28}
\]

Například pro \(Z=-s^2\), \(s>0\), je levá strana
\(-2s^{-1}\arctan(r/s)\); identita \(\arctan(r/s)+\arctan(s/r)=\pi/2\)
dává (BF28) se \(\sqrt Z=is\). Analytické pokračování pak dokazuje právě
vnější větev (BF4). Pro \(|Z|\le r^2/4\) je \(|J_r|\le4/r\).

Tím je (BF3) přesná identita s

\[
B_{\mathrm{loc}}(K,Z)=\frac i\pi
\left[A(K,Z)J_r(Z)-\int_{-r}^rH_*(x^2,Z,K)\,dx\right]
+\int_{-r}^r h_-R_A P_x\frac{dx}{2\pi}.
\tag{BF29}
\]

Všechny členy na pravé straně jsou holomorfní v obyčejném Z.

\hypertarget{vuxfdslovnuxe9-meze-koeficientux16f-a-skuteux10dnuxfd-jednostrannuxfd-dux16fsledek}{%
\subsection{Výslovné meze koeficientů a skutečný jednostranný
důsledek}\label{vuxfdslovnuxe9-meze-koeficientux16f-a-skuteux10dnuxfd-jednostrannuxfd-dux16fsledek}}

Na polydisku (BF24) lze použít \(|h_-|<5\), \(|r_b|<64\), \(|P_x|<8\).
Proto \(|a|<2^{12}\). Funkce A je tudíž holomorfní a \(|A|<2^{12}\) na
\(|K-K_0|<2^{-60}\), \(|Z|<2^{-124}\), což je menší než čtverec
přípustného x-poloměru. Na polovičních vnitřních oborech Cauchy dává
bezpečně

\[
|A_K|\le2^{73},\qquad |A_Z|\le2^{137}.
\]

Z (BS31),(BF8),(BF24) je \(A(K_0,0)>3/256\). Na (BF2) se tato hodnota
změní o nejvýše \(2^{73}2^{-88}+2^{137}2^{-152}=2^{-14}\). Tedy
\(|A|>2^{-8}\), jak bylo tvrzeno. Pro reálné K,Z je A reálná a kladná.

Pro meze B použijme Cauchyho koeficientovou mez
\(|a_{2j}|\le2^{12}(2^{-61})^{-2j}\). Jelikož \(r=2^{-61}/8\), pro
\(|x|\le r\), \(|Z|\le r^2/4\) platí

\[
|H_*|\le2\,2^{12}(2^{-61})^{-2}=2^{135},
\quad
\left|\int_{-r}^rH_*\,dx\right|\le2^{72},
\] \[
|A J_r|\le2^{78}.
\]

Regulární integrál v (BF29) je omezen číslem
\(2r\cdot5\cdot2^{34}\cdot8<2^{-23}\), kde zanedbání \(1/(2\pi)<1\)
zvětšuje mez. Součet těchto tří mezí je menší než \(2^{80}\). Tím je
celé číselné tvrzení (BF2),(BF3) dokázáno.

Pro reálné K a \(0<Z<2^{-152}\) je navíc

\[
\boxed{\operatorname{Re}F_{2,\mathrm{loc}}(K,E)
=\frac{A(K,Z)}{\sqrt Z}.}
\tag{BF30}
\]

Zde je podstatné, že (BF11),(BF18) drží \textbf{celý místní objektový
interval v jeho fyzikální mezeře}, v níž má f\_A mimo svůj atom čistě
imaginární okrajovou hodnotu na reálné ose. Reziduum \(r_b\), pulzní
váha h a \(P_x\) jsou reálné. Po odečtení \(2ir_b/(\omega-\omega_b)\) je
tedy také \(R_A\) pro reálné argumenty čistě imaginární. V (BF29) jsou
A, \(J_r\), \(H_*\) reálné, takže celý \(B_{\mathrm{loc}}\) je čistě
imaginární. Oba atomové průchody \(x=\pm\sqrt Z\) jsou již zahrnuty v
(BF28); jejich kladný reálný příspěvek je přesně (BF30).

Zbytek původní úplné konvoluce má nezápornou reálnou spektrální hustotu:
je integrálem úplných Carathéodoryho spektrálních měr s původními
nezápornými pulzními vahami. Odříznutý P-interval je při reálném K
skutečný kladný díl této míry. Proto

\[
\operatorname{Re}F_2(K,E)\ge
\operatorname{Re}F_{2,\mathrm{loc}}(K,E)
\ge\frac{2^{-8}}{\sqrt Z}.
\tag{BF31}
\]

Nerovnost se chápe pro obvyklé okrajové hodnoty, případně téměř všude; v
nekonečné hustotě má přenosový sloupec příslušný mezní význam. Zde se
nepotřebuje číselná horní mez komplementu.

Na původním reálném energetickém okně je již dokázáno
\(\operatorname{Re}F_1\ge1/16\), \(|F_1|\le2\). Protože
\(\operatorname{Re}F_2\ge0\), pro konečné nenulové \(F_2\) platí

\[
|4+F_1F_2|=|F_2|\,|F_1+4/F_2|
\ge |F_2|\operatorname{Re}(F_1+4/F_2)
\ge |F_2|/16\ge 2^{-12}/\sqrt Z.
\]

Přesná algebraická identita

\[
V-\binom{-2/F_1}{0}
=\frac{(8/F_1,4)^T}{4+F_1F_2}
\tag{BF32}
\]

a \(8/|F_1|+4\le132<2^8\) dokazují (BF6). Jde o odhad \textbf{celého V},
nikoli o dosazení neúplného \(F_{2,\mathrm{loc}}\) do jeho jmenovatele.

\hypertarget{co-pux159esnux11b-zbuxfdvuxe1-pro-oboustrannuxfd-normuxe1lnuxed-tvar-celuxe9ho-v}{%
\subsection{Co přesně zbývá pro oboustranný normální tvar celého
V}\label{co-pux159esnux11b-zbuxfdvuxe1-pro-oboustrannuxfd-normuxe1lnuxed-tvar-celuxe9ho-v}}

Položme \(F_{2,\mathrm{comp}}=F_2-F_{2,\mathrm{loc}}\), se stejnými
pohyblivými koncovými body, takže tato rovnost zachovává úplný integrál.
Úplné \(F_1\) má podle (FQ9) holomorfní pokračování se
\(\operatorname{Re}F_1>1/16\), \(|F_1|<3\) na \(|E-E_0|<2^{-30}\). Celé
naše okolí (BF2) v něm leží, neboť podle (BF20)

\[
|E-E_0|\le8|K-K_0|+|Z|<2^{-84}<2^{-30}.
\]

Jestliže se nyní nezávisle doloží, že na společném okolí je komplement
holomorfní v K,Z, pišme \(B=B_{\mathrm{loc}}+F_{2,\mathrm{comp}}\). S
\(s=\sqrt Z\) pak dostaneme přesně

\[
\boxed{
V(K,E_b(K)+s^2)
=\frac{(-2A-2sB,\,4s)^T}{F_1A+s(4+F_1B)}.}
\tag{BF33}
\]

Funkce A,B na pravé straně jsou vyhodnoceny v \((K,s^2)\); úplné
\(F_1\), které na K samostatně nezávisí, v \(E_b(K)+s^2\). Nenulový
součin \(F_1(E_b(K))A(K,0)\) ruší inverzní odmocninovou singularitu
celého přenosu. Zůstává holomorfní funkce s, s rozvojem

\[
V(K,E_b+s^2)=\binom{-2/F_1(E_b)}{0}
+s\binom{8/(F_1(E_b)^2 A(K,0))}
{4/(F_1(E_b)A(K,0))}+O(s^2).
\tag{BF34}
\]

Koeficient s je nenulový; druhá složka jej dokazuje bez možnosti
zrušení. Závislost \(F_1\) na E přidává až členy řádu s². Hlavní
odmocninový koeficient přenosu nezávisí na hodnotě B.

Pro zcela číselný oboustranný poloměr zbývá dodat na společném
komplexním polydisku mez pro B. Ve vzorci níže ponecháme obecná označení
pro tři meze

\[
|B|\le M_B,\qquad c_1\le|F_1|\le M_1,
\qquad c_1>0,
\tag{BF35}
\]

z nichž \(c_1=1/16\), \(M_1=3\) jsou již certifikovány pomocí (FQ9).
Zbývající \(M_B\) se musí vztahovat na celý B, nejen na
\(B_{\mathrm{loc}}\).

S již dokázaným \(a_*=2^{-8}\) postačuje

\[
|s|\le\frac{c_1a_*}{2(4+M_1M_B)};
\]

jmenovatel (BF33) pak má absolutní hodnotu nejméně \(c_1a_*/2\), a
přesná obdoba (BF32) dává

\[
\left\|V-\binom{-2/F_1(E_b+s^2)}{0}\right\|_2
\le |s|\frac{2(8/c_1+4)}{c_1a_*}.
\tag{BF36}
\]

Konkrétně dodatečná mez \(|B|\le2^b\), \(b\ge0\), na společném oboru
(BF2) by spolu s již doloženým (FQ9) dala

\[
\boxed{
\begin{gathered}
|K-K_0|<2^{-88},\qquad
|s|\le\min\{2^{-77},\,2^{-(b+16)}\}\\
\Longrightarrow\quad
\left\|V-\binom{-2/F_1(E_b+s^2)}{0}\right\|_2
\le2^{21}|s|.
\end{gathered}}
\tag{BF37}
\]

První poloměr ve společném minimu zajišťuje
\(|Z|=|s|^2\le2^{-154}<2^{-152}\), tedy přísnou podmínku (BF2), i když
je b malé. Skutečně \(c_1a_*=2^{-12}\),
\(2(4+3\cdot2^b)\le14\cdot2^b<2^{b+4}\) a
\(2(8/c_1+4)/(c_1a_*)=264\cdot2^{12}<2^{21}\). Jde o výslovnou
implikaci; dodatečná číselná mez komplementu se zde nepředstírá.

Oddělení od běžných kontinuálních záhybů o více než 1/300 v K vylučuje
shodu obou již klasifikovaných stacionárních tříd. Samo však nepřiděluje
číselnou mez celé zbývající periodě. Je třeba společně ošetřit její
nestacionární póly, objektové hrany, společné příspěvky pólu s hranou a
pulzní šev podle původního fyzikálního obvodu. Mez \(2^{80}\) pro
\(B_{\mathrm{loc}}\) nelze použít jako mez pro toto B bez dalšího
důkazu.

Tento místní výsledek nepřidává nový globální čas rozptylu: ostré konce
původního energetického okna a ostatní fyzikální prahy mají stále svůj
vlastní úplný chybový účet. Platné starší časové meze se nemění. Nově
jsou doloženy místní kvadratické souřadnice, správně normalizovaný
nenulový odmocninový koeficient úplného místního objektového příspěvku a
nepodmíněný jednostranný odhad (BF6) pro skutečný celý přenos.


\section{Úplný časový účet a přesně zbývající práce}

\subsection{Které nevyčíslené kroky byly odstraněny}

QA, QB a QC dávají pro tři dříve pouze kvalitativně
kontrolované typy místních příspěvků konkrétní obvody,
jejich vnější fyzikální identifikaci a všechny parametrické
derivace. U pólu s hranou je společný obvod nutný pro
zachování rušení. U kužele obě hrany odcházejí stejným
směrem. U výměny pulzu se přesný původní součet nejprve
analyticky rozštěpí; dva zachované členy poté používají
opačné obvody. Tento rozdíl je doložen před vnějším limitem.

FQ kontroluje \textbf{celý} první kontakt. BF dává dva
Rouchého kroky, číselnou Morseovu souřadnici, přesný
integrál s původním faktorem \(i/\pi\) a meze úplného
místního Laurentova zbytku. Na reálné spektrální straně
z kladnosti celého komplementu plyne BF6 přímo pro celý
přenosový sloupec. Nevyžaduje neznámou horní mez komplementu.

\subsection{Nejmenší doplnění pro oboustranný místní přenos}

V okolí každého vázaného prahu BF přesně vyčleňuje jeden
interval integrační hybnosti. Definujme
\[
F_{2,\mathrm{comp}}=F_2-F_{2,\mathrm{loc}}
\]
se stejnými pohyblivými konci. Pro oboustranný analytický
tvar celého \(V\) zbývá doložit na explicitním společném
okolí holomorfii tohoto \textbf{úplného fyzikálního}
komplementu a konkrétní mez
\[
|F_{2,\mathrm{comp}}|\le M_{\rm comp}<\infty.
\]
Má-li se použít menší okolí než BF2, musí být jeho poloměr
rovněž výslovně určen. Holomorfie na neurčeném dostatečně
malém okolí k číselnému času nestačí.

S \(M_B=2^{80}+M_{\rm comp}\) potom stačí
\[
|\sqrt Z|<\min\left\{2^{-76},\,
\frac{2^{-13}}{4+3M_B}\right\}.
\]
Původní jmenovatel po vynásobení \(\sqrt Z\) je oddělen
od nuly číslem \(2^{-13}\). Z BF33--BF37 plyne plný
oboustranný odhad s koeficientem \(2^{21}\).
Číslo \(2^{80}\) již doložené v BF omezuje pouze
\(B_{\mathrm{loc}}\). Není horní mezí celého \(B\).

Tento potřebný krok je konkrétní úloha pro již zadané
funkce. Z dosavadní práce neplyne nutnost nového fyzikálního
předpokladu. Rovněž se netvrdí, že holomorfie komplementu
nebo uvedená číselná mez už byla dokázána.

\subsection{Další kroky před menším globálním časem}

I po doplnění komplementu u vázaných prahů zbývají dvě
skutečné běžné kontinuální stacionární křivky. Je třeba
vyčíslit jejich úplné místní periody a jejich přenos na
\(V\), dále pokrýt regulární části původního kruhu a
sečíst odhady v původní normě. Číselný rozestup CS12
usnadňuje oddělení těchto úloh, ale žádnou z nich sám
neprovádí.

Následný Fourierův odhad musí zahrnout celý kruh \(K\),
oba ostré konce původního \(I\), úplný sloupec \(V\)
a konečný přijímač. Jednostranné BF6 samo neurčuje
globální Sobolevovu konstantu ani součet časového chvostu.
Proto z něj tento dodatek neodvozuje menší \(N_+\).

\subsection{Kontrolované vyhlazení zachovává původní zdroj}

Ostré energetické konce lze v důkazu dočasně ošetřit
hladkým skalárem \(0\le q_\eta\le1\), který je roven
jedné na \([E_*-h+\eta,E_*+h-\eta]\) a má podporu
uvnitř \(I\). Pomocný vektor se nepřenormovává.
Původní rovnoměrná energetická míra dává přesně
\[
\nu^2=\|(1-q_\eta(H_0))\psi\|^2
=\frac1{2h}\int_I|1-q_\eta(E)|^2\,\mathrm dE
\le\frac\eta h.
\]
Vrácení pomocného vektoru k původnímu vstupu ve srovnání
dvou izometrických vývojů stojí nejvýše
\[
2\nu\le2\sqrt{\eta/h}.
\]
Pro \(\eta=10^{-18}\) je tato chyba menší než
\(6.4\cdot10^{-8}\). Projektor \(Q\) a absolutní
pravděpodobnost se tím nepředefinují. Pro časový závěr
je ještě třeba zaplatit derivace hladkého přechodu,
rostoucí s \(\eta^{-1}\), a sečíst skutečný chvost.

\subsection{Dosavadní samostatný konečný čas zůstává beze změny}

Předchozí interně přezkoumaný obecný návrh [S4], za
převzatých operátorových a rozptylových předpokladů,
uvádí konečné počty původních diskrétních makrokroků
\[
N_-=10^{16},\qquad N_+=2^{81\,000\,000\,000\,000},
\]
a úplnou časovou chybu včetně přijímače
\[
\varepsilon_-(N_-)+\varepsilon_+(N_+)
<2.310002\cdot10^{-7}<10^{-6}.
\]
S původními \(\varepsilon_S=\varepsilon_D=10^{-6}\)
zůstává jeho konzervativní amplitudový účet
\[
\left|\sqrt{p_{\rm finite}}-\sqrt{P_0}\right|
<2.24\cdot10^{-6}.
\]
Původní absolutní dolní mez
\(L=293166031669590845829/2^{96}>(60\cdot10^{-6})^2\)
proto v témže matematickém modelu dává
\[
p_{\rm finite}>(57.76\cdot10^{-6})^2
=3.3362176\cdot10^{-9}.
\]
Tyto časy a tato pravděpodobnost jsou převzaty z [S4].
Nový místní dodatek je nezmenšuje, ani nově nepočítá
absolutní pravděpodobnost. [S4] je samostatný analytický
návrh této relace, nikoli veřejně přijatý časový důkaz
v PR~\#1817. Konečný matematický protokol není
uskutečněným laboratorním experimentem. Převod makrokroku
na fyzikální sekundy v původním modelu doložen není.

\section{Veřejná návaznost, způsob kontroly a zdroje}

\subsection{Kontrola před tímto pokračováním}

Kontrola pouze čtením proběhla 29.~10.~2026
v~10:53:18--10:53:25 CET (Europe/Rome, UTC+1).
PR~\#1817 a PR~\#1815 zůstávaly otevřenými nesloučenými
návrhy na commitech
\nolinkurl{66ef8220df3ad5764ca941b84223a9f46cf5411f} a
\nolinkurl{32306016908c0097650277aa416fc0152b2f81a4}.
Jejich seznamy komentářů a review byly prázdné.
Rezervace \#1816 zůstávala otevřená, nezměněná, bez
komentářů. Kolekce issues a PR aktualizovaných po
10:19 CET byla prázdná. Nové veřejné pokračování
časového důkazu při této kontrole nalezeno nebylo.

Závěrečná kontrola 29.~10.~2026 v~10:18:13 UTC
(11:18:13 CET) znovu ověřila PR~\#1817, rezervaci
\#1816 a hlavní větev. Od 09:53:25 UTC nebyla
aktualizována žádná issue ani PR; uvedené stavy a
commity zůstaly shodné.

Hlava hlavní větve byla
\nolinkurl{a193e389fb1ddb0bdc89165b015c95a17eb1e425}.
STATUS nadále označoval Public Canon v106 jako ACTIVE,
s obsahovým commitem
\nolinkurl{9242db4f040d67d9417eacd06612230b74248cff}.
Tento stav není tvrzením o novém dokončeném úplném
replay. Nevydává se nový veřejný výsledek ani se
nemění obsah Canonu.

\subsection{Rozsah vnitřního přezkumu}

Nové důkazy byly nezávisle analyticky přezkoumány
v této pracovní relaci. Kontrola pokryla nenulové
meze původních jmenovatelů na obou listech, počet
kořenů obou Rouchého kroků, vnější znaménka,
periodicitu úplného prvního kontaktu, správnou
pulzní maticovou identitu a původní normalizaci
odmocninového koeficientu. Přezkum rozlišuje místní
\(B_{\mathrm{loc}}\) a dosud neomezený úplný komplement.

Jde o ruční analytické odhady s racionálními rezervami.
Nové numerické hledání, CAS vyhodnocení ani vědecký
výpočet nebyly provedeny. Sazba a kontrola dokumentu
nejsou vydávány za numerickou certifikaci matematických
tvrzení. Tento přezkum není externím přijetím ani
součástí dřívějšího auditu 23 skupin a 130 podmínek.

Práce je samostatná, \textbf{NON-CANONICAL / L1}.
Rezervace \#1816 ani veřejné návrhy nebyly měněny.

\section*{Odkazy na neměnné podklady}

\begin{description}[leftmargin=9mm,style=nextline]
\item[S1] Původní absolutní zachycení, část I, zejména
oddíly 1, 4--8 a 11:
\url{https://github.com/mathorn1973/twist-j/blob/32306016908c0097650277aa416fc0152b2f81a4/notes/C-BOUND-PAIR-ABSOLUTE-CAPTURE-N/PROOF.md}.
\item[S2] Úplný algebraický audit původního čtyřkanálového
operátoru, zejména oddíly 9--10:
\url{https://github.com/mathorn1973/twist-j/blob/66ef8220df3ad5764ca941b84223a9f46cf5411f/notes/C-BOUND-PAIR-RESOLVENT-AUDIT-N/PROOF.md}.
\item[S3] Konečná příprava, čtení a původní časová mezera:
\url{https://github.com/mathorn1973/twist-j/blob/66ef8220df3ad5764ca941b84223a9f46cf5411f/notes/C-BOUND-PAIR-RESOLVENT-AUDIT-N/FINITE_PROTOCOL.md}.
\item[S4] Samostatný předchozí návrh této relace,
\textit{Konečný čas rozptylu původního paketu TWIST-J},
29.~10.~2026, 33 stran. Zde zejména obecný konečný
časový účet a oddíl 14 s rovnicemi PT:
\url{twist_j_konecny_cas.tex}.
\item[S5] Samostatný předchozí dodatek,
\textit{Vázané větve a skutečné rozptylové prahy TWIST-J},
29.~10.~2026, 28 stran. Úplný soupis pólů, rovnice
BS25--BS31, CS12 a společný kvalitativní obvod PE:
\url{twist_j_vazane_vetve.tex}.
\item[S6] Veřejné návrhy a oddělená rezervace:
\url{https://github.com/mathorn1973/twist-j/pull/1817};
\url{https://github.com/mathorn1973/twist-j/pull/1815};
\url{https://github.com/mathorn1973/twist-j/issues/1816}.
\item[S7] Provozní označení aktivního Canonu:
\url{https://github.com/mathorn1973/twist-j/blob/a193e389fb1ddb0bdc89165b015c95a17eb1e425/STATUS.md}.
\end{description}

\end{document}